\documentclass[11pt]{article}
\usepackage[utf8]{inputenc}
\usepackage[T1]{fontenc}
\usepackage{amsmath,amssymb,amsthm}
\usepackage[margin=1.1in]{geometry}
\usepackage{microtype}
\usepackage{url}
\usepackage[hidelinks]{hyperref}

\theoremstyle{plain}
\newtheorem{theorem}{Theorem}[section]
\newtheorem{lemma}[theorem]{Lemma}
\newtheorem{proposition}[theorem]{Proposition}
\newtheorem{corollary}[theorem]{Corollary}
\theoremstyle{definition}
\newtheorem{definition}[theorem]{Definition}
\newtheorem{example}[theorem]{Example}
\newtheorem{remark}[theorem]{Remark}

\newcommand{\Zhat}{\widehat{\mathbb{Z}}}
\newcommand{\N}{\mathbb{N}}
\newcommand{\dlog}{\delta_{\log}}
\newcommand{\cU}{\widehat{U}}
\newcommand{\Ahat}{\widehat{A}}
\newcommand{\Kh}{\widehat{K}}

\title{A profinite reduction of Erd\H{o}s Problem 25}
\author{Ibrahim Mian \qquad Shayaan Siddique\\[2pt]
\normalsize Millennium Research\\[2pt]
\normalsize \texttt{\{ibby,shayaan\}@millenniumresearch.ai}\\
\normalsize \texttt{ibrahimnmian@gmail.com}, \texttt{shayaansiddique02@gmail.com}}
\date{}

\begin{document}
\maketitle

\begin{abstract}
Let $1\le n_1<n_2<\cdots$ be integers with residues $a_i$, and
let $A$ be the set of integers $n$ such that for every $i$ either
$n<n_i$ or $n\not\equiv a_i\pmod{n_i}$.  Erd\H{o}s asked whether
$A$ must have a logarithmic density.  We recast the question in
the profinite completion $\Zhat$, where each class is a ball of
Haar measure $1/n_i$ and the survivors have a Haar measure $\nu$
that always exists, and prove unconditionally that the upper
logarithmic density of $A$ is at most $\nu$, so that the reverse
inequality would settle the problem.  We prove the reverse inequality in
a series of regimes: convergent drift;
pairwise incompatible, pairwise coprime, aligned and boundedly
misaligned systems; systems admitting a cylinder decomposition;
and, through reduction theorems -- exhaustion of dying
cylinders, confinement to protectors of summable measure, and a
theorem for coprime-layered trees of cylinders proved by a
martingale argument -- systems of divergent drift, positive
survivor measure and unbounded misalignment whose classes are
organised along nested protectors, pairwise coprime protectors,
or protectors that are the pairwise products of a set of primes,
each under explicit and restrictive hypotheses on the layers.
A second-moment condition, with a lemma restoring the zero--one
law for families whose indices do not escape to infinity, gives
a sufficient condition for $\nu=0$.  We exhibit the step of the
Davenport--Erd\H{o}s argument that shifted residues break, refute
four natural routes by explicit construction, bound a fifth, and show that every
system is majorised by a saturated one whose classes are the
maximal cylinders meeting its survivor set in a null set.  What remains is stated as
a checklist that a counterexample must pass in full
and, without congruences, as a question about a closed subset of
$\Zhat$ of positive measure whose maximal holes have divergent
total measure.  The problem is not resolved here.
\end{abstract}

{\small\noindent\textit{2020 Mathematics Subject Classification.}
11B05, 11N35.

\noindent\textit{Keywords.} Sets of multiples, logarithmic density,
congruence systems, profinite integers, Davenport--Erd\H{o}s.}


\section{Introduction}\label{sec:intro}

Let $1\le n_1<n_2<\cdots$ be a strictly increasing sequence of
integers and let $a_i\in\mathbb{Z}$ be arbitrary.  The $i$th
class kills the set
\[
  K_i=\{x\ge n_i:\ x\equiv a_i\ (\mathrm{mod}\ n_i)\},
\]
and the survivor set is $A=\N\setminus\bigcup_i K_i$.  We write
$\delta(E)$ for the natural density of a set $E$ when it exists,
and $\dlog(E)$ for its logarithmic density.  Erd\H{o}s
asked whether $A$ must possess a logarithmic density,
\[
  \dlog(A)=\lim_{X\to\infty}\frac{1}{\log X}\sum_{m\in A,\ m<X}\frac1m .
\]
The truncation $x\ge n_i$ is essential: without it each class is
a full residue class and the question becomes a statement about
sets of multiples only in the aligned case.

When every $a_i\equiv0$ the killed set is a set of multiples and
the answer is affirmative by the theorem of Davenport and
Erd\H{o}s \cite{DE1936,DE1951}; see \cite{Hall} for a modern
account.  The general problem has resisted precisely because
shifting the residues destroys the divisibility structure that
their proof consumes.  Our purpose is to say exactly what that
destruction costs, to prove as much of the statement as survives
it, and to isolate what remains.

We do not claim the individual ingredients are deep.  The
profinite upper bound is elementary; the regime results rest on
the Chinese remainder theorem, Davenport--Erd\H{o}s and a
standard Tauberian argument.  What we offer is an
accounting: the step of the classical proof that misalignment
breaks, a weighted form of Davenport--Erd\H{o}s that survives
it, the regimes and reductions under which the reverse
inequality is recovered, a single uniform statement whose truth would settle the
rest, and four natural approaches shown by
explicit construction not to work, together with a bound on a
fifth, the most natural attempt at a counterexample.

\subsection*{Results}
Sections~\ref{sec:profinite} and~\ref{sec:upper} set up the
profinite framework and prove the unconditional half of the
problem (Theorem~\ref{thm:upper}), together with the threshold
convention and the restriction lemma by which systems pass to
cylinders.  Section~\ref{sec:lemmas} records two lemmas on sets
of multiples (Lemmas~\ref{lem:tail} and~\ref{lem:weight}), the
second a weighted form of Davenport--Erd\H{o}s.
Section~\ref{sec:regimes} proves the reverse inequality for
convergent drift, incompatible systems, translated systems and
finitely many aligned channels (Theorems~\ref{thm:conv},
\ref{thm:H}, \ref{thm:transl}, \ref{thm:chan}), proves $\nu=0$
under coprime divergence (Theorem~\ref{thm:D}), and gives the
second-moment criterion for $\nu=0$ (Theorem~\ref{thm:quasi},
Lemma~\ref{lem:escape}).  Section~\ref{sec:destop} exhibits the
step of the Davenport--Erd\H{o}s argument that shifted residues
break (Proposition~\ref{prop:break}).  Section~\ref{sec:cyl}
proves what survives it: bounded misalignment and cylinder
decompositions (Theorems~\ref{thm:bounded}, \ref{thm:cyl},
\ref{thm:cylapp}).  Section~\ref{sec:reductions} contains the
reduction theorems: exhaustion and saturation
(Theorems~\ref{thm:exhaust}, \ref{thm:sat}), summable protectors
(Theorem~\ref{thm:confined}), coprime-layered trees
(Theorem~\ref{thm:tree}), and coprime and pairwise-product
protectors (Theorems~\ref{thm:coprot}, \ref{thm:overlap}).
Section~\ref{sec:reduction} gives a uniform statement that
suffices for the rest (Theorem~\ref{thm:reduction}), its
excursion form, and the checklist and closed-set question in
which we leave the problem.  Section~\ref{sec:routes} refutes
four routes by construction and bounds a fifth;
Section~\ref{sec:verification} describes the verification.  We emphasise at the
outset that the problem is \emph{not} resolved here.

\medskip\noindent\textbf{The problem in the literature.}
The question is Problem~25 on the Erd\H{o}s problems
site~\cite{Bloom}, listed as open, and as
a special case of Problem~486, in which several residues may be
excluded modulo each $n$; its statement has been formalised in
Lean.  Tenenbaum has recorded that Erd\H{o}s was much interested
in whether the Davenport--Erd\H{o}s theorem extends to non-zero
residues, and that he himself has thought about it without
finding a promising starting point~\cite{TaoThread}.  A
neighbouring question, Problem~281, asks: if $n_1<n_2<\cdots$ is
such that for every choice of residues the integers avoiding all
the classes have density $0$, must there be, for every
$\varepsilon>0$, a $k$ such that for every choice of residues the
integers avoiding the first $k$ classes have density below
$\varepsilon$?  It was answered affirmatively in 2026 by the observation
(recorded in~\cite{TaoRogers}) that an old theorem of Rogers
(Halberstam--Roth~\cite{HaRo66}, Ch.~V) applies: given
$\varepsilon$, the 1936 Davenport--Erd\H{o}s theorem gives a $k$
for which the \emph{aligned} survivors of the first $k$ classes
have density below $\varepsilon$, and Rogers' inequality then
bounds the survivors of the same $k$ moduli with any residues by
that aligned density.  It does not bear on the existence of the
density.  The
only prior progress on Problem~25 itself that we know of is a
2026 note of Chojecki~\cite{Choj26}, which proves that the
natural density exists when the drift converges (his
Theorem~3.1; Theorem~\ref{thm:conv} below is its logarithmic
form) and when all the moduli are pairwise coprime (his
Theorem~3.2; for divergent drift this is the case of
Theorem~\ref{thm:D} below in which the whole family is coprime),
and reduces the
general problem to a uniform harmonic estimate for the finite
``quotient sieves'' by which each class is first reached, proving
two obstructions: the tail union admits no vanishing logarithmic
bound, and a prime-power tower gadget gives a single first-kill
set a transient harmonic contribution far larger than its
eventual density loss.  In the language of this paper the first
is the failure of the union bound under divergent drift; the
second is the transient-spike phenomenon whose aggregate effect
on the logarithmic average Proposition~\ref{prop:spike} bounds.  Apart from those two
theorems we have found no prior source for the results below;
the weighted Davenport--Erd\H{o}s lemma in particular is an
elementary variant of the classical Dirichlet-series proof and
we claim no more for it than its use here.  We note also
where the neighbouring literature stops.  The lower bounds on
the density of integers left uncovered by a system of congruences
-- the inequality of Filaseta, Ford, Konyagin, Pomerance and
Yu~\cite{FFKPY}, a product term less a pairwise-dependence
correction, and Hough's relative form of the local
lemma~\cite{Hough} -- concern the density of the periodic set
uncovered by finitely many classes, that is the profinite
measure $\nu$ of the present paper.  They bound $\nu$ from
below; the open half of Problem~25 is a lower bound on the
integer survivors at scales far below the period, and no
estimate of that kind appears in this literature.  Erd\H{o}s's
own criterion~\cite{Er48} for a set of multiples to have a
natural density -- that integers whose least divisor in the set
lies in $(n^{1-\varepsilon},n]$ be rare -- concerns the same
top-of-range phenomenon of which Proposition~\ref{prop:spike}
controls only the single dyadic window $(X/2,X]$, and only in
logarithmic density; the range $(X^{1-\varepsilon},X/2]$ is
exactly what remains uncontrolled here.  The theory of Behrend sequences -- sets of
multiples of full density -- is the aligned case's version of
the question when $\nu=0$, and it is instructive how hard even
that is.  Hall and Tenenbaum~\cite{HT92} and
Tenenbaum~\cite{Ten96} obtain sharp criteria only for block
sequences, in the form of Borel--Cantelli conditions whose exponent
\[
  \alpha=\frac{1-\log2}{1-(1+\log\log2)/\log2}\approx3.565,
\]
in Tenenbaum's words, measures the level of mutual dependence
among the events (\cite{Ten96}, the display following (1.7)); the method conditions on the product $n_k$ of an
integer's prime factors below $\exp\exp k$ and iterates bounds
on the conditional probability of a new divisor, which is the
multiplicative counterpart of the cylinder filtration in
Theorem~\ref{thm:tree}.  Tenenbaum describes effective general
criteria as very difficult, if not hopeless.  For shifted
classes, sufficient conditions for $\nu=0$ are given here by
the cruder second-moment route of Theorems~\ref{thm:D}
and~\ref{thm:quasi} and Lemma~\ref{lem:escape}; the open half of
Problem~25 concerns $\nu>0$, a case which for sets of multiples
is entirely Davenport--Erd\H{o}s's and has no counterpart in
that literature.

\section{The profinite framework}\label{sec:profinite}

Write $\Zhat=\varprojlim \mathbb{Z}/m\mathbb{Z}$ for the
profinite integers with normalised Haar measure $\mu$.

\begin{definition}
To the class $(n,a)$ associate the ball
$B(n,a)=a+n\Zhat\subseteq\Zhat$, which is clopen with
$\mu(B(n,a))=1/n$.  Write $B_i=B(n_i,a_i)$ and
\[
  \cU=\bigcup_i B_i,\qquad \Ahat=\Zhat\setminus\cU,\qquad \nu=\mu(\Ahat).
\]
Throughout, a subscript on $\nu$ names the subsystem whose
survivor measure is meant -- the first $N$ classes ($\nu_N$),
the classes of modulus at most $X$ ($\nu_X$), a set of classes
or groups ($\nu_T$, $\nu_J$) -- and a prime marks a reduced
system; for a cylinder, $\nu_c$ or $\nu_C$ is the survivor
measure of the induced system there, that is, the measure
relative to the cylinder.  We fix once and for all an enumeration of the
classes by non-decreasing modulus, ties in any fixed order;
``the first $N$ classes'' refers to it.  Nothing below depends on
the choice: finite heads enter only through limits over $N$, and
the same quantities can be defined by a modulus cutoff instead.
A system is a \emph{set} of classes: no ball is listed twice,
a repetition killing nothing new.  Several classes may share a
modulus, a slight widening of the problem in which Erd\H{o}s's
case is that of strictly increasing moduli; the reductions of
Section~\ref{sec:reductions} produce systems of this wider kind.  Since a modulus
$n$ carries at most $n$ residues, a system has at most $n$
classes of modulus $n$, so bounded modulus still means finitely
many classes, and an infinite system has unbounded moduli.  We reserve the hat $\cU$ for this subset of $\Zhat$
and write $U=\bigcup_i K_i\subseteq\N$ for the set of killed
integers; likewise $\Ahat\subseteq\Zhat$ is the profinite
survivor set and $A\subseteq\N$ the integer one.  The profinite
and integer objects are never identified.
\end{definition}

Since $\cU$ is open, $\mu(\cU)$ exists, and therefore
\emph{$\nu$ always exists}, with no hypothesis on the system.  This is the
first gain of the profinite viewpoint: the candidate value of
the answer is available before anything has been proved.

\begin{lemma}\label{lem:compat}
Two classes $(n,a)$ and $(n',a')$ have a common solution if and
only if $a\equiv a'\pmod{\gcd(n,n')}$, and this holds if and
only if $B(n,a)\cap B(n',a')\ne\emptyset$.  A family of classes
is pairwise compatible if and only if it has a simultaneous
solution \emph{in $\Zhat$}, equivalently if and only if the
corresponding balls have a common point.  For infinite families
that solution need not be an integer; see
Remark~\ref{rem:notaligned}.
\end{lemma}

\begin{proof}
The first assertion is the two-modulus case of the Chinese
remainder theorem.  For the second, note that
$\Zhat\cong\prod_p\mathbb{Z}_p$ and a ball is the product of its
$p$-adic components.  In $\mathbb{Z}_p$ balls are ultrametric:
if two intersect then one contains the other.  Hence any finite
pairwise-intersecting family has a ball of smallest radius,
which is contained in all the others, so the family has the
finite intersection property; compactness of $\mathbb{Z}_p$
then gives a nonempty intersection in each coordinate.
\end{proof}

\begin{definition}
A \emph{channel} is a pairwise compatible family of classes; by
Lemma~\ref{lem:compat} this is exactly a family of balls sharing
a common point of $\Zhat$.  A channel is \emph{aligned} if that
common point may be taken in $\mathbb{Z}$, i.e.\ if there is an
integer $c$ with $a_i\equiv c\pmod{n_i}$ for all its members.
The \emph{drift} of a family is $\sum 1/n_i$.
\end{definition}

\begin{remark}\label{rem:notaligned}
The two notions coincide for finite families, where the Chinese
remainder theorem produces an integer solution modulo the lcm,
but not in general.  An infinite channel has a common point of
$\Zhat$, and $\Zhat$ is uncountable while $\mathbb{Z}$ is
countable, so a coherent sequence of residues modulo $2^k$ may
be chosen avoiding every integer; the resulting family is
pairwise compatible but aligned to no integer.  The distinction
is not cosmetic: the translation argument of
Theorem~\ref{thm:transl} requires an integer, so it applies to
aligned channels only.
\end{remark}

\section{The unconditional half}\label{sec:upper}

\begin{theorem}\label{thm:upper}
For every system, $\overline{\delta}_{\log}(A)\le\nu$.  No
hypothesis on the moduli or residues is required.
\end{theorem}

\begin{proof}
Fix $N$ and let $A_N$ be the survivor set of the first $N$
classes.  The truncated condition $x\equiv a_i$ and $x\ge n_i$
differs from the profinite condition $x\in B_i$ only at
$x=a_i\bmod n_i$ when that representative is smaller than $n_i$,
so the two survivor sets of the first $N$ classes differ by at
most $N$ integers.  A finite set has density zero, and the
profinite survivor set of finitely many classes is a finite
union of residue classes, hence periodic; so $A_N$ has a natural
density and a logarithmic density, both equal to
$\nu_N:=\mu(\Zhat\setminus\bigcup_{i\le N}B_i)$.  Since
$A\subseteq A_N$ we get $\overline{\delta}_{\log}(A)\le\nu_N$
for every $N$, and $\nu_N\downarrow\nu$ by continuity of measure
from above.
\end{proof}

\begin{corollary}\label{cor:nuzero}
If $\nu=0$ then $\dlog(A)$ exists and equals $0$.
\end{corollary}

\begin{proof}
By Theorem~\ref{thm:upper}, $\overline{\delta}_{\log}(A)\le\nu=0$,
and logarithmic densities are non-negative.
\end{proof}

\medskip\noindent\textbf{Thresholds.}
Passing to a cylinder, as we shall do repeatedly, changes the
point at which a class begins to kill.  A \emph{threshold
system} consists of classes $(n_i,a_i;t_i)$, the class killing
the integers $x\ge t_i$ with $x\equiv a_i\pmod{n_i}$, where
$1\le t_i\le n_i$ and, for some constant $M'$ and all but
finitely many $i$, $t_i\ge n_i/M'$.  The systems of the
introduction are the case $t_i=n_i$.  Thresholds are invisible
to every profinite quantity, since balls do not see them.  On
the integer side they enter in exactly five ways.  A finite
head remains periodic up to finitely many integers.  The
harmonic sum of a class over $[1,X]$ is at most
$(\log X+1)/n_i+1/t_i\le(\log X+1+M')/n_i$, so every tail
bound of the form $\sum_{\text{tail}}(\log X+O(1))/n_i$ keeps
its shape with a larger constant.  And a class whose induced
condition is a pure divisibility $m\mid y$ is unchanged by a
threshold at most $m$, since the only multiple of $m$ below $m$
is $0$.  Finally, a class of modulus $n$ lying inside a cylinder
of modulus $d$ kills only integers $x\ge t_i\ge n/M'\ge d/M'$,
so its kills lie in $C\cap[d/M',\infty)$, whose harmonic sum up
to $X$ is at most $(\log X+\log M'+M')/d$: the tail bounds of
Section~\ref{sec:reductions}, which sum such quantities over cylinders, keep their
shape with a larger constant.  There is one further point of
contact, on the side of natural rather than logarithmic density:
a class of modulus $n_i$ kills at most $X/n_i+1$ integers below
$X$, the term $1$ being absent for a standard class, and only
classes with $t_i\le X$, hence $n_i\le M'X$, contribute; when
$\sum_i1/n_i<\infty$ the number of such classes is $o(X)$, since
$\#\{i:Y/2<n_i\le Y\}\le Y\sum_{n_i>Y/2}1/n_i=o(Y)$ summed over
dyadic ranges, so the accumulated extra terms are $o(X)$.  Every
result of this paper is stated for standard systems and holds
for threshold systems with the same proof, thresholds entering
only through these five facts: Theorems~\ref{thm:upper}
and~\ref{thm:D} through the first, Theorem~\ref{thm:conv} and
Lemma~\ref{lem:tail} through the second,
Theorem~\ref{thm:bounded} through the third,
Theorems~\ref{thm:exhaust}, \ref{thm:confined} and
Lemma~\ref{lem:light} through the fourth, with the system's
constant $M'$ in place of $1$, and Theorem~\ref{thm:H} through
the fifth -- with one exception, Theorem~\ref{thm:sat}, which is
stated for standard systems and is false for threshold systems:
the classes $(2^k,2^{k-1};2^{k-1})$, $k\ge1$, kill every positive
integer, while their balls fill $2\Zhat$ up to the point $0$, so
$\nu=\tfrac12$ and the saturated system $\{2\Zhat\}$ has integer
survivors of logarithmic density $\tfrac12$.  Its proof uses
that a class kills nothing below its modulus.  The theorem is
applied only to the original system.  Induced systems arise in two
ways only.  Theorems~\ref{thm:bounded}, \ref{thm:cyl}
and~\ref{thm:cylapp} restrict to the cylinders of one fixed
modulus $M$, and Lemma~\ref{lem:restrict} gives those induced
systems the constant $2MM'$; the tree theorem restricts to the
finitely many cells of one fixed depth, each with its own such
constant.  Theorem~\ref{thm:exhaust} also speaks of the induced
system of each group on its own cylinder, of unbounded modulus,
but only to apply Corollary~\ref{cor:nuzero} to each one
separately, which needs no threshold constant at all.  No
argument restricts repeatedly with unbounded moduli: the
reduction theorems of Section~\ref{sec:reductions} act on the
original system, and the fourth fact bounds their tail kills by
the original thresholds.
The correspondence between a system and an induced one holds up
to the finitely many integers below the cylinder's modulus,
which no density sees.

\begin{lemma}[Restriction]\label{lem:restrict}
Let $c+M\mathbb{Z}$ be a cylinder with $0\le c<M$ and write
$x=c+My$, $y\ge1$.  A class
$(n,a;t)$ either misses the cylinder, when
$\gcd(n,M)\nmid a-c$, or becomes in $y$ the class of modulus
$m=n/\gcd(n,M)$, with the residue determined by $a$ and $c$, and
threshold $t'=\max\bigl(1,\lceil(t-c)/M\rceil\bigr)\le m$.  For
all but finitely many classes $t'\ge m/(2MM')$, so the induced
family is a threshold system, and, apart from the finitely
many integers below $M$, the integers of the cylinder surviving
the original system are exactly $c+My$ for $y$ surviving the
induced one.
\end{lemma}

\begin{proof}
The congruence $c+My\equiv a\pmod n$ has a solution iff
$\gcd(n,M)\mid a-c$, and then its solutions form one class
modulo $m$.  The condition $c+My\ge t$ is $y\ge(t-c)/M$, and
$(t-c)/M\le n/M\le m$.  Since $t\ge n/M'$ for all but finitely
many classes, $t'\ge n/(MM')-1\ge m/(MM')-1\ge m/(2MM')$ once
$m\ge2MM'$, which excludes only the finitely many classes with
$n\le 2M^2M'$ (finitely many, a system having at most $n$
classes of modulus $n$).  The last sentence is the definition of the
induced system.
\end{proof}

\begin{remark}[Rogers' theorem]\label{rem:rogers}
The measure $\nu$ is largest when all residues vanish.  A
theorem communicated by Rogers and published only in
Halberstam and Roth~\cite[Ch.~V]{HaRo66} -- see Tao's
exposition~\cite{TaoRogers} -- states that in a finite cyclic
group the union of cosets $a_j+H_j$ is at least as large as the
union of the subgroups $H_j$.  Applied in $\mathbb{Z}/L_N\mathbb{Z}$
to the first $N$ classes and letting $N\to\infty$, it gives
$\nu\le\nu_{\mathrm{al}}$, where $\nu_{\mathrm{al}}$ is the profinite survivor measure
of the aligned system, which by Davenport--Erd\H{o}s is the
logarithmic density of the integers divisible by no modulus.
With Theorem~\ref{thm:upper} the chain
$\overline{\delta}_{\log}(A)\le\nu\le\nu_{\mathrm{al}}=\dlog\{m:\ n_i\nmid m
\ \forall i\}$ shows that the survivors of \emph{any} residue
choice have upper logarithmic density at most the logarithmic
density of the non-multiples of the same moduli.
The inequality runs the wrong way for the open half of the
problem, which is a lower bound on the survivors.
\end{remark}

So every system whose survivors vanish in measure is settled by
the upper bound alone, with no further hypothesis; the open
problem lives entirely in the range $\nu>0$.

A sufficient form of Erd\H{o}s's question is therefore the
reverse inequality
$\underline{\delta}_{\log}(A)\ge\nu$, which asserts that the
classes beyond any finite head destroy no more than their own
measure.  It is sufficient and not equivalent: a system whose
survivors had a logarithmic density strictly below $\nu$ would
answer the question affirmatively while violating it, and
nothing here excludes that.  Every positive result below proves
the inequality, and the reductions of Section~\ref{sec:reduction}
are stated for it.

\section{Two lemmas on sets of multiples}\label{sec:lemmas}

We isolate two facts used in Sections~\ref{sec:regimes}--\ref{sec:reductions}.  Write $M(D)$
for the set of positive multiples of members of $D$, and
$D^{(N)}$ for the $N$ smallest members of $D$.

\begin{lemma}[Vanishing tails]\label{lem:tail}
Let $D$ be a set of moduli and let $D_1\subseteq D_2\subseteq
\cdots$ be any increasing exhaustion of $D$ by finite sets.
Then $\dlog\bigl(M(D)\setminus M(D_N)\bigr)\to0$ as
$N\to\infty$.  In particular this holds for $D_N=D^{(N)}$, the
$N$ smallest members.
\end{lemma}

\begin{proof}
Each $M(D_N)$ is a finite union of sets of multiples, hence
periodic, so it has a natural density $\delta(M(D_N))$, and
these increase with $N$ because the sets do.  The family
$(D_N)$ is cofinal in the directed family of finite subsets of
$D$: every finite $F\subseteq D$ lies in some $D_N$, since
$F$ is finite and the $D_N$ exhaust $D$.  Therefore
\[
  \lim_N\delta\bigl(M(D_N)\bigr)
  \ =\ \sup_{F\subseteq D\ \text{finite}}\delta\bigl(M(F)\bigr)
  \ =\ \dlog\bigl(M(D)\bigr),
\]
the last equality being the theorem of Davenport and Erd\H{o}s
\cite{DE1936,DE1951} (see \cite[Ch.~0]{Hall}), which in
particular gives that $M(D)$ has a logarithmic density.  Note
that the limit is thus the same for every exhaustion, not only
for the one by the $N$ smallest members.

The union
\[
  M(D)=M(D_N)\ \sqcup\ \bigl(M(D)\setminus M(D_N)\bigr)
\]
is disjoint, and logarithmic density is finitely additive on
disjoint sets whenever two of the three quantities exist; here
$M(D)$ and $M(D_N)$ both have logarithmic densities, so the
difference set has one, equal to
$\dlog(M(D))-\delta(M(D_N))$, which tends to $0$.
\end{proof}

\begin{lemma}[Weighted Davenport--Erd\H{o}s]\label{lem:weight}
Let $\varphi:\N\to[0,1]$ satisfy $\varphi(d)\le\varphi(n)$
whenever $d\mid n$.  Then
$\lim_{X\to\infty}\frac{1}{\log X}\sum_{n\le X}\varphi(n)/n$
exists.
\end{lemma}

\begin{proof}
If $\varphi\equiv0$ the assertion is trivial, so assume
$\varphi(n_0)>0$ for some $n_0$; then
$F(\sigma):=\sum_{n\ge1}\varphi(n)n^{-\sigma}>0$ for
$\sigma>1$, and division by $F$ below is legitimate.  The series
converges absolutely for $\sigma>1$ since $\varphi\le1$.
Monotonicity gives
\[
  \sum_{d\mid n}\varphi(d)\Lambda(n/d)
  \ \le\ \varphi(n)\sum_{d\mid n}\Lambda(n/d)
  \ =\ \varphi(n)\log n,
\]
the inequality that drives the Davenport--Erd\H{o}s argument and
which reappears as \eqref{eq:lemma1} below.
Multiplying by $n^{-\sigma}$ and summing, the left side is the
Dirichlet convolution $F(\sigma)\cdot(-\zeta'/\zeta)(\sigma)$
and the right side is $-F'(\sigma)$, so
\[
  -F'(\sigma)\ \ge\ F(\sigma)\Bigl(-\frac{\zeta'}{\zeta}\Bigr)
  (\sigma),
  \qquad\text{i.e.}\qquad
  \frac{d}{d\sigma}\log\frac{F(\sigma)}{\zeta(\sigma)}\ \le\ 0 .
\]
Hence $F/\zeta$ is non-increasing on $(1,\infty)$ and takes
values in $[0,1]$, so $L=\lim_{\sigma\to1^+}F(\sigma)/\zeta(\sigma)$
exists.  Since $(\sigma-1)\zeta(\sigma)\to1$ we get
$(\sigma-1)F(\sigma)\to L$, and as the coefficients $\varphi(n)$
are non-negative, Karamata's Tauberian theorem for the
Laplace--Stieltjes transform of the non-decreasing function
$t\mapsto\sum_{\log n\le t}\varphi(n)/n$ (see
\cite[Ch.~II.7]{Tenenbaum}) converts $(\sigma-1)F(\sigma)\to L$
into $\sum_{n\le X}\varphi(n)/n\sim L\log X$ when $L>0$, the
stated logarithmic mean.  When $L=0$ the conclusion
$\sum_{n\le X}\varphi(n)/n=o(\log X)$ follows from positivity
alone: with $\alpha(t)=\sum_{\log n\le t}\varphi(n)/n$ one has
$\alpha(t)\le e\,F(1+1/t)$, since every term with $\log n\le t$
carries the factor $n^{-1/t}\ge e^{-1}$ in $F(1+1/t)$, so
$\alpha(t)/t\le e\,t^{-1}F(1+1/t)\to0$.  Note that
nothing beyond $0\le\varphi\le1$ and divisibility-monotonicity
was used; in particular $\varphi$ need not be an indicator
function, which is the case of Davenport--Erd\H{o}s.
\end{proof}

\section{Five regimes in which the reverse inequality holds}\label{sec:regimes}

\begin{theorem}[Convergent drift]\label{thm:conv}
If $\sum_i 1/n_i<\infty$ then $\dlog(A)$ exists and equals
$\nu$, for every choice of residues.
\end{theorem}

\noindent(This is Theorem~3.1 of Chojecki~\cite{Choj26}, who
obtains the natural density; the proof below is included because
its harmonic union bound is used throughout.)

\begin{proof}
We first record the harmonic union bound, which is what the
logarithmic setting requires: a bound on natural density would
not suffice, since a set of small natural density can carry
large logarithmic mass.

The elements of $K_i$ below $X$ are at least $n_i$ and lie in an
arithmetic progression of modulus $n_i$, so the $j$th of them is
at least $j\,n_i$; hence
\[
  \frac{1}{\log X}\sum_{x\in K_i,\ x\le X}\frac1x
  \ \le\ \frac{1}{n_i\log X}\sum_{j\le X/n_i}\frac1j
  \ \le\ \frac{1}{n_i}\Bigl(1+\frac{1}{\log X}\Bigr).
\]
Summing and using $1_{\bigcup K_i}\le\sum 1_{K_i}$ pointwise,
\[
  \limsup_{X\to\infty}\frac{1}{\log X}
  \sum_{x\in\bigcup_{i>N}K_i,\ x\le X}\frac1x
  \ \le\ \sum_{i>N}\frac1{n_i}.
\]

Now given $\varepsilon>0$ choose $N$ with
$\sum_{i>N}1/n_i<\varepsilon$.  Then $A\subseteq A_N$ and
$A_N\setminus A\subseteq\bigcup_{i>N}K_i$, whose upper
logarithmic density is below $\varepsilon$ by the display above.
Since $A_N$ is periodic apart from a finite symmetric
difference with the profinite survivor set of the first $N$
classes, $\dlog(A_N)=\nu_N$; hence the upper and lower
logarithmic
densities of $A$ satisfy
$\nu_N-\varepsilon\le\underline{\delta}_{\log}(A)
\le\overline{\delta}_{\log}(A)\le\nu_N$, so they differ by at
most $\varepsilon$.  As $\varepsilon$ was
arbitrary they agree; the common value is at most $\nu$ by
Theorem~\ref{thm:upper} and at least $\nu_N-\varepsilon\ge
\nu-\varepsilon$ for every $\varepsilon$, so it is $\nu$.
\end{proof}

\begin{lemma}[Disjointness bound]\label{lem:A}
If the classes are pairwise incompatible then
$\sum_i 1/n_i\le 1$, and the constant $1$ is attained.
\end{lemma}

\begin{proof}
Incompatible classes have no common solution, so the sets $K_i$
are pairwise disjoint; each has density $1/n_i$ since truncation
removes finitely many elements.  For any finite subfamily the
disjoint union lies in $\N$, giving $\sum 1/n_i\le1$; take the
supremum.  Attainment: the class $(2,0)$ together with $(12,r)$
for the six odd $r$ is pairwise incompatible and has drift
$\tfrac12+6\cdot\tfrac1{12}=1$.  (Before truncation these
residue classes cover $\mathbb{Z}$; after truncation the odd
integers below $12$ survive, which does not affect the drift
computation.)
\end{proof}

\begin{theorem}[Incompatible systems]\label{thm:H}
If the classes are pairwise incompatible then the natural
density of $A$ exists and equals $1-\sum_i 1/n_i$; in
particular $\dlog(A)$ exists and equals $\nu$.
\end{theorem}

\begin{proof}
The $K_i$ are pairwise disjoint, so for each $N$ the density of
$\bigcup_{i\le N}K_i$ is exactly $\sum_{i\le N}1/n_i$, giving
$\underline{\delta}\bigl(\bigcup_i K_i\bigr)\ge\sum_{i\le N}1/n_i$
for every $N$.  For the reverse bound, the elements of $K_i$
below $X$ are at least $n_i$ and lie in a progression of modulus
$n_i$, so there are at most $X/n_i$ of them (at most $X/n_i+1$
for a threshold class, the extra terms summing to $o(X)$ by the
fifth threshold fact); disjointness then
gives $\bigl|\bigcup_i K_i\cap[1,X]\bigr|\le X\sum_i 1/n_i+o(X)$,
whence $\overline{\delta}\le\sum_i 1/n_i$.  By
Lemma~\ref{lem:A} that sum converges, so upper and lower
densities agree.
\end{proof}

\begin{theorem}[Coprime divergence]\label{thm:D}
If some pairwise coprime subfamily has divergent drift then
$A$ has density $0$, and $\dlog(A)=0=\nu$.
\end{theorem}

\noindent(When the whole family is pairwise coprime this is
contained in Theorem~3.2 of Chojecki~\cite{Choj26}.)

\begin{proof}
Coprime moduli give independent kill events by the Chinese
remainder theorem, so the survivor density through the first
$N$ members of the subfamily is exactly $\prod(1-1/n_i)$, which
tends to $0$ when $\sum 1/n_i$ diverges.  Above the largest of
those $N$ moduli, $A$ is contained in each of these periodic
sets, so $\delta(A)=0$.  The same product tends to
$0$ in measure, so $\nu=0$.
\end{proof}

The coprime hypothesis can be replaced by a second-moment
condition on the kills, grouped by smallest prime.

\begin{theorem}[Quasi-independence]\label{thm:quasi}
For each prime $p$ let $\Kh_p\subseteq\Zhat$ be the union of the
balls of the classes whose modulus has smallest prime factor
$p$.  If $\sum_p\mu(\Kh_p)=\infty$ and
\[
\limsup_{N\to\infty}\
\frac{\sum_{p,p'\le N}\mu(\Kh_p\cap \Kh_{p'})}
     {\bigl(\sum_{p\le N}\mu(\Kh_p)\bigr)^2}<\infty,
\]
then $\nu=0$, and hence $\dlog(A)=0$.
\end{theorem}

\begin{proof}
Haar measure on $\Zhat=\prod_\ell\mathbb{Z}_\ell$ is the product
of the Haar measures of the coordinates, so the coordinates are
independent.  Every prime dividing a modulus with smallest prime
factor $p$ is at least $p$, so $\Kh_p$ is determined by the
coordinates at primes $\ell\ge p$.  Hence
$E=\{x:\ x\in \Kh_p\text{ for infinitely many }p\}$ lies in the
tail $\sigma$-algebra of the coordinates, and by Kolmogorov's
zero--one law $\mu(E)\in\{0,1\}$.  The Kochen--Stone
lemma~\cite{KS}, the second-moment form of Borel--Cantelli,
gives under the two hypotheses
$\mu(E)\ge1/\limsup_N(\cdots)>0$.  So $\mu(E)=1$.  Every point
of $E$ lies in a ball, so the survivor set in $\Zhat$ is
contained in $\Zhat\setminus E$ and $\nu=0$;
Corollary~\ref{cor:nuzero} finishes.
\end{proof}

The hypothesis that $\Kh_p$ be indexed by smallest primes is
only a way of forcing the sets to move off to infinity in the
coordinates.  What the proof uses is the following, and it will
be needed in a case where the index is not a prime at all.

\begin{lemma}[Escaping families]\label{lem:escape}
Let $(\Kh_j)$ be one countable family of measurable subsets of
$\Zhat$, and for each $P$ let $\mathcal{K}_P$ be the set of its
members determined by the coordinates at primes exceeding $P$;
these sets decrease as $P$ grows.  Suppose
that for every $P$ the members of $\mathcal{K}_P$ can be
enumerated so that their measures have divergent sum and the
second-moment ratios along some subsequence of that
enumeration have limsup at most a constant $\rho$ independent of
$P$.  Then almost every point of $\Zhat$
lies in infinitely many $\Kh_j$.
\end{lemma}

\begin{proof}
Fix $P$ and let $E_P$ be the set of points lying in infinitely
many members of $\mathcal{K}_P$.  Kochen--Stone~\cite{KS} gives
$\mu(E_P)\ge1/\rho$.  The families $\mathcal{K}_P$ decrease as
$P$ grows, so the $E_P$ decrease, and $\mu(\bigcap_PE_P)
=\lim_P\mu(E_P)\ge1/\rho>0$.  Each $E_P$ is determined by the
coordinates above $P$, so the intersection lies in the tail
$\sigma$-algebra and has measure $1$ by Kolmogorov's zero--one
law.  It is contained in the set of points lying in infinitely
many $\Kh_j$.
\end{proof}

\begin{remark}
The distinction matters.  In Theorem~\ref{thm:quasi} the sets
are indexed by their smallest prime and escape automatically.
In Theorem~\ref{thm:overlap} below they are indexed by pairs of
primes, and infinitely many of them involve any given small
prime, so the set of points covered infinitely often is
\emph{not} a tail event; only after discarding the pairs that
touch the first $\pi(P)$ primes does the argument apply, and it
is the limit over $P$ that restores the zero--one law.
\end{remark}

\begin{remark}
Theorem~\ref{thm:D} is the case in which the subfamily's
smallest primes are distinct and each $\Kh_p$ is a single ball:
then $\mu(\Kh_p\cap \Kh_{p'})=\mu(\Kh_p)\mu(\Kh_{p'})$ by the Chinese
remainder theorem and the ratio tends to $1$.  Read in
contrapositive the theorem is a structure statement: a system
with $\nu>0$ has either $\sum_p\mu(\Kh_p)<\infty$ -- the kills
through each smallest prime are summable, as when the smallest
primes of the moduli are themselves sparse -- or smallest-prime
kills so positively correlated that the ratio is unbounded.
The many-protector systems of Section~\ref{sec:reductions} --
systems whose classes fall into groups each confined to a
cylinder, its \emph{protector}; when the protector's modulus is a
prime $p_j$ we also call $p_j$ the protector -- fall on either
side according to the sparsity of the protectors: a class $p_jm$ has
smallest prime $\min(p_j,p(m))$, so $\Kh_p$ collects only the
protectors $p_j\ge p$ and $\sum_p\mu(\Kh_p)\approx\sum_j\pi(p_j)/p_j$,
which behaves like $\sum_j1/\log p_j$: convergent for
$p_j\ge\exp(j^{2})$, say, and divergent for exponentially or
polynomially sparse $p_j$, where positive survivor measure then
forces the correlated branch.  Divergent drift alone gives
neither.
\end{remark}

\begin{theorem}[Translation]\label{thm:transl}
If all classes satisfy $a_i\equiv c\pmod{n_i}$ for one fixed
integer $c\ge0$, then $A=c+\{y:\ n_i\nmid y\ \text{for all }i\}$
up to a finite set, where $y$ runs over the positive integers,
and $\dlog(A)$ exists and equals $\nu$.
\end{theorem}

\begin{proof}
For $x>c$ put $y=x-c\ge1$.  The class $(n_i,c\bmod n_i)$ kills
$x$ iff $n_i\mid y$ and $x\ge n_i$; and if $n_i\mid y$ then
$y\ge n_i$, so $x\ge n_i$ holds automatically, the threshold
mattering only at $x=c$, which is excluded.  So apart from the
finitely many $x\le c$, $x\in A$
iff no $n_i$ divides $y$.  The set of such $y$ has a logarithmic
density by Davenport--Erd\H{o}s, equal to the limit of the
densities of its finite heads, which is $\nu$; translation by
$c$ and a finite change preserve logarithmic density.
\end{proof}

\begin{remark}
The hypothesis $c\ge0$ is necessary and the finiteness depends
on it.  The exceptions are the integers lying below their
class's truncation point, that is the representatives of $c$
modulo each $n_i$.  For $c\ge0$ and $n_i>c$ that representative
is $c$ itself, so at most $c+1$ distinct integers arise however
many classes there are.  For $c<0$ this fails badly: with
$c=-1$ the representative modulo $n$ is $n-1<n$, so \emph{every}
class contributes a fresh exception and the symmetric difference
is infinite.  Note also that an infinite aligned channel
determines $c$ uniquely, since two aligning integers would
differ by a multiple of arbitrarily large moduli; so $c$ cannot
be normalised to be nonnegative.  Theorem~\ref{thm:chan} is
therefore proved below without appealing to this theorem.
\end{remark}

\begin{theorem}[Finitely many aligned channels]\label{thm:chan}
If the classes are a union of finitely many \emph{aligned}
channels then $\dlog(A)$ exists and equals $\nu$.
\end{theorem}

\begin{proof}
Let $C_1,\dots,C_r$ be the channels, $c_j$ the integer aligning
$C_j$, and $D_j$ the set of moduli occurring in $C_j$.  We argue
without any sign hypothesis on $c_j$, and in particular without
invoking Theorem~\ref{thm:transl}.

Fix $N$ and let $D_j^{(N)}$ consist of the
$\min\{N,|D_j|\}$ smallest members of $D_j$, so that
$D_j^{(N)}=D_j$ once $N\ge|D_j|$; channels with
$D_j^{(N)}=\varnothing$ contribute nothing and are omitted from
the maximum defining $B(N)$ below.  Let $S_N$ be the survivor
set of the finitely many classes of the channels whose moduli
lie in the respective $D_j^{(N)}$, each channel being a set of
classes and the classes assigned to it exhausted by modulus.  Being a finite union
of truncated residue classes, $S_N$ is periodic apart from
finitely many integers and has a density $d_N$
equal to its logarithmic density.

We claim that for every $x$ exceeding
$B(N):=\max_j\bigl(\max D_j^{(N)}+|c_j|\bigr)$,
\[
  x\in S_N\setminus A
  \ \Longrightarrow\
  x-c_j\in M(D_j)\setminus M\bigl(D_j^{(N)}\bigr)
  \ \text{for some } j,
\]
where $M(D)$ is the set of positive multiples of members of $D$.
Indeed such an $x$ is killed by some class $(n,a)$ of some $C_j$
with $n\notin D_j^{(N)}$, so $n\mid x-c_j$, and $x>|c_j|$ makes
$x-c_j$ positive, whence $x-c_j\in M(D_j)$.  And $x$ survives
every head class, so for $x$ larger than every head modulus we
have $m\nmid x-c_j$ for all $m\in D_j^{(N)}$, that is
$x-c_j\notin M(D_j^{(N)})$.

By Lemma~\ref{lem:tail},
$\dlog\bigl(M(D_j)\setminus M(D_j^{(N)})\bigr)\to0$ as
$N\to\infty$, and logarithmic density is invariant under
translation by a fixed integer of either sign, up to discarding
finitely many initial terms -- which the bound $B(N)$ above has
already excluded.  Given
$\varepsilon>0$ choose $N$ so that these $r$ quantities total
less than $\varepsilon$.  Since $A\subseteq S_N$ and
$S_N\setminus A$ is contained, apart from the at most $B(N)$
integers below the bound, in
$\bigcup_j\bigl(c_j+(M(D_j)\setminus M(D_j^{(N)}))\bigr)$, the
upper and lower logarithmic densities of $A$ each lie within
$\varepsilon$ of $d_N$.  As $\varepsilon$ was arbitrary,
$\dlog(A)$ exists; and since $d_N=\nu_N\downarrow\nu$ while the
tail bound tends to $0$, its value is $\nu$.
\end{proof}

\begin{remark}
The hypothesis that the channels be \emph{aligned} cannot be
dropped from the proof: by Remark~\ref{rem:notaligned} an
infinite channel need not admit an integer translate, and
without one Theorem~\ref{thm:transl} does not apply.  Whether
the conclusion survives for merely compatible channels we do not
know; such a channel has unbounded misaligned parts, beyond
the reach of Theorem~\ref{thm:bounded}, and none of the later
reductions is tailored to it.
\end{remark}

Two structural facts produce compatible channels, which are
aligned when finite and so feed Theorem~\ref{thm:chan} in that
case.  Pairwise
coprime classes are automatically compatible, since $\gcd=1$
makes the residue condition vacuous.  And if a subfamily has
$\gcd(n_i,n_j)=d$ for every pair, then partitioning it by
residue modulo $d$ yields at most $d$ channels, one for each
residue class that actually occurs, so divergent drift for the
family forces divergent drift for one of them, carrying at least
$1/d$ of the total.

\section{Where Davenport--Erd\H{o}s stops}\label{sec:destop}

The engine of the Davenport--Erd\H{o}s argument
\cite{DE1936} is the inequality
\begin{equation}\label{eq:lemma1}
  f(n)\log n\ \ge\ \sum_{d\mid n} f(d)\Lambda(n/d),
\end{equation}
where $f$ is the indicator of the killed set.  Its proof uses
one property and nothing else: the killed set is closed upward
under divisibility, so $f(d)\le f(n)$ whenever $d\mid n$, and
the right side is at most $f(n)\sum_{d\mid n}\Lambda(n/d)=f(n)\log n$.

\begin{proposition}\label{prop:break}
Shifted classes destroy that closure, and \eqref{eq:lemma1}
fails.  Consider the finite system with moduli
$2\le n\le 39$ and $a_n=1$.  Then $31$ is killed while $62$ is
not, so the killed set is not closed upward under divisibility;
and \eqref{eq:lemma1} fails at $n=42$, where the left side is
$0$ while the divisors $6$, $14$ and $21$ are all killed, giving
a right side of $\log 2+\log 3+\log 7>3.7$.
\end{proposition}

\begin{remark}
The restriction to a finite head is not a defect of the example
but a feature of the setting.  In the full system $(n,1)$,
$n\ge2$, every $x\ge3$ is killed by the class $n=x-1$, so
\eqref{eq:lemma1} holds trivially with equality.  The
Davenport--Erd\H{o}s argument is applied to the finite heads
$U_k$ and passed to the limit, so it is exactly at the finite
heads that \eqref{eq:lemma1} must hold, and exactly there that
misalignment breaks it.
\end{remark}

\section{Bounded misalignment and cylinder decompositions}\label{sec:cyl}

What survives is a splitting.

\begin{definition}
For a class $(n,a)$ write $n=uw$ with $\gcd(u,w)=1$ and
$u\mid a$; call $u$ the \emph{aligned part} and the least such
$w$ the \emph{misaligned part}.  Then
\[
  x\equiv a\ (\mathrm{mod}\ n)
  \iff u\mid x \ \text{ and }\ x\equiv a\ (\mathrm{mod}\ w).
\]
\end{definition}

\begin{theorem}\label{thm:bounded}
If the misaligned parts all divide a fixed $M$, then $\dlog(A)$
exists and equals $\nu$.
\end{theorem}

\begin{proof}
Fix a residue $c$ modulo $M$ and write $C_c=\{x\equiv c\ (M)\}$.
Since $w_i\mid M$, every $x\in C_c$ satisfies
$x\equiv c\pmod{w_i}$, so class $i$ can kill some element of
$C_c$ only when $a_i\equiv c\pmod{w_i}$: the cylinder
\emph{selects} which classes act on it.  For the selected
classes the remaining condition is $u_i\mid x$.  Writing
\[
  D_c=\{u_i:\ a_i\equiv c\ (\mathrm{mod}\ w_i)\},
\]
the killed set inside $C_c$ is therefore $M(D_c)\cap C_c$, up to
the truncation.

Let $U=\bigcup_i K_i$ be the killed integers, $U_r$ the union
of the first $r$ of them, and
$B_r=\max_{i\le r}n_i$.  If $x\in C_c$ lies in $U\setminus U_r$
and $x>B_r$, then some class kills it, so $x\in M(D_c)$; and no
class of the head kills it, so, $x$ exceeding every head
modulus, $u_j\nmid x$ for every $u_j\in D_{c,r}$, where
$D_{c,r}$ collects the selected $u_j$ with $j\le r$.  Hence
\[
  (U\setminus U_r)\cap C_c\ \subseteq\
  \bigl(M(D_c)\setminus M(D_{c,r})\bigr)\cap C_c
\]
apart from the finitely many $x\le B_r$.

Intersecting with $C_c$ only shrinks a set, so the upper
logarithmic density of the left-hand side is at most that of
$M(D_c)\setminus M(D_{c,r})$, which tends to $0$ as
$r\to\infty$ by Lemma~\ref{lem:tail}, the sets $D_{c,r}$
exhausting $D_c$.  Summing the $M$ cylinders,
$\overline{\delta}_{\log}(U\setminus U_r)\to0$.

Now given $\varepsilon>0$ choose $r$ with
$\overline{\delta}_{\log}(U\setminus U_r)<\varepsilon$.  The set
$U_r$ is a finite union of truncated progressions, hence
periodic apart from finitely many integers, so it has a density
$1-\nu_r$.  As $U_r\subseteq U\subseteq U_r\cup(U\setminus U_r)$,
the upper and lower logarithmic densities of $U$ each lie within
$\varepsilon$ of $1-\nu_r$, so they agree.  Hence $\dlog(U)$
exists and equals $\lim_r(1-\nu_r)=1-\nu$, so $\dlog(A)$ exists
and equals $\nu$.
\end{proof}

\begin{remark}\label{rem:gap}
Two more obvious routes to this theorem do not work, and we
record them because each looks plausible.  One cannot finish by
applying Davenport--Erd\H{o}s cylinderwise: a set of multiples
intersected with a fixed residue class modulo $M$ is not a set
of multiples, and substituting $x=c+My$ turns divisibility into
a congruence.  Nor does the divisibility-monotone weight
\[
  \varphi(x)=\tfrac1M\#\{c\bmod M:\ u_i\mid x\ \text{and}\
  a_i\equiv c\ (\mathrm{mod}\ w_i)\ \text{for some }i\}
\]
suffice, though Lemma~\ref{lem:weight} does give it a
logarithmic mean: $\varphi$ records the cylinders in which $x$
\emph{would} be killed, whereas $x$ inhabits one cylinder, and
the two quantities differ.  For the classes $(6,2)$, $(10,4)$,
$(14,6)$, $(15,5)$, $(21,7)$ with $M=105$, both limits are
exact: the killed set is periodic modulo $210$ with density
$73/210=0.3476$, which is the limit of $F$, the harmonic mass of
the killed set defined in Section~\ref{sec:reduction}; while
$\varphi$ is periodic modulo $70$ and its logarithmic mean is
$169/490=0.3449$.  What makes the proof above work is that the
cylinder \emph{selects} the acting classes, which is exactly
what $w_i\mid M$ provides.
\end{remark}

The hypothesis can be relaxed, and the relaxation already
reaches systems of unbounded misalignment.

\begin{corollary}\label{cor:mixed}
Suppose the classes split as $S_1\cup S_2$, where every
misaligned part occurring in $S_1$ divides a fixed $M$, and
$\sum_{i\in S_2}1/n_i<\infty$.  Then $\dlog(A)$ exists and
equals $\nu$.  The misaligned parts of the system as a whole
may be unbounded.
\end{corollary}

\begin{proof}
Given $\varepsilon>0$ choose $N$ with
$\sum_{i\in S_2,\ i>N}1/n_i<\varepsilon$ and let $T$ consist of
$S_1$ together with the first $N$ members of $S_2$.  Every
misaligned part occurring in $T$ divides
$M_1=\operatorname{lcm}(M,w_{i_1},\dots,w_{i_N})$, which is
finite, so Theorem~\ref{thm:bounded} applies to $T$ and
$\dlog(\operatorname{surv}T)=\nu_T$.  Now
$A\subseteq\operatorname{surv}T$ and
$\operatorname{surv}T\setminus A$ is contained in the union of
the remaining members of $S_2$, whose upper logarithmic density
is at most $\sum_{i\in S_2,\ i>N}1/n_i<\varepsilon$ by the
harmonic bound of Theorem~\ref{thm:conv}.  Hence the upper and
lower logarithmic densities of $A$ each lie within $\varepsilon$
of $\nu_T$, and likewise $|\nu_T-\nu|\le\varepsilon$; let
$\varepsilon\to0$.
\end{proof}

So unboundedness alone is not the obstruction: it must be
unboundedness carrying divergent drift.  For the system $(n,1)$
the misaligned parts are
$2,3,\dots,n,\dots$, no finite $M$ exists, and the cylinder
decomposition is infinite.

A second extension reaches systems that Theorem~\ref{thm:bounded}
and Corollary~\ref{cor:mixed} both miss.  On the cylinder
$x\equiv c\pmod M$, write $x=c+My$; by Lemma~\ref{lem:restrict}
the classes meeting the cylinder become a threshold system in
$y$, the \emph{induced system} at $c$, whose survivors are
exactly the indices of the integers of the cylinder surviving
the original.

\begin{theorem}[Cylinder decomposition]\label{thm:cyl}
Suppose there is a modulus $M$ such that for every residue $c$
modulo $M$ the induced system at $c$ either has profinite
survivor measure $0$, or has bounded misaligned parts.  Then
$\dlog(A)$ exists and equals $\nu$.
\end{theorem}

\begin{proof}
The sets $A\cap C_c$ for $c$ modulo $M$ are disjoint with union
$A$, and there are finitely many, so it suffices that each has a
logarithmic density.  Fix $c$ and let $S$ be the survivor set of
the induced system at $c$, so that $A\cap C_c=\{c+My:\ y\in S\}$.
Then
\[
  \frac{1}{\log X}\sum_{t\in A\cap C_c,\ t\le X}\frac1t
  =\frac{1}{\log X}\sum_{y\in S,\ y\le (X-c)/M}\frac{1}{c+My}
  \ \longrightarrow\ \frac{1}{M}\,\dlog(S),
\]
since $1/(c+My)=\bigl(1+O(1/y)\bigr)/(My)$ and
$\log\bigl((X-c)/M\bigr)/\log X\to1$.  By hypothesis $\dlog(S)$
exists.  In the first case the induced system has profinite
survivor measure $0$, so $\overline{\delta}_{\log}(S)\le0$ by
Theorem~\ref{thm:upper} applied to that system, giving
$\dlog(S)=0$.  In the second case Theorem~\ref{thm:bounded} applies and gives
$\dlog(S)=\nu_c$, the survivor measure of the induced system.
Summing over $c$, and using $\sum_c\nu_c/M=\nu$, gives the
result.
\end{proof}

The hypothesis asks more than the argument needs: it is enough
that almost every cylinder be resolved.

\begin{theorem}[Approximate cylinder decomposition]\label{thm:cylapp}
Suppose that for every $\varepsilon>0$ there is a modulus $M$
such that the residues $c$ modulo $M$ whose induced system is
neither of profinite survivor measure $0$ nor of bounded
misaligned parts number fewer than $\varepsilon M$.  Then $\dlog(A)$ exists and equals $\nu$.
\end{theorem}

\begin{proof}
Fix $\varepsilon$ and such an $M$, and split the residues into
the resolved set $G$ and the rest $B$, with $|B|<\varepsilon M$.
Write $C_G$ and $C_B$ for the corresponding unions of
cylinders, so that $A=(A\cap C_G)\sqcup(A\cap C_B)$.  As in
Theorem~\ref{thm:cyl}, $A\cap C_G$ is a finite disjoint union of
sets each having a logarithmic density, so it has one; call it
$g$.  And $A\cap C_B\subseteq C_B$, a union of $|B|$ residue
classes modulo $M$, of logarithmic density $|B|/M<\varepsilon$.
Hence
\[
  g\ \le\ \underline{\delta}_{\log}(A)
   \ \le\ \overline{\delta}_{\log}(A)\ \le\ g+\varepsilon .
\]
The value $g$ depends on $\varepsilon$, but the gap bound does
not; as $\varepsilon$ was arbitrary the upper and lower
logarithmic densities agree, and since $g$ is the survivor
measure of the resolved cylinders, which lies within
$\varepsilon$ of $\nu$, the common value is $\nu$.
\end{proof}

\begin{remark}
Theorem~\ref{thm:cyl} is the case $\varepsilon=0$.  The
relaxation is reachable where the strict form is not: a class
of modulus $n$ meets $M/\gcd(n,M)$ of the $M$ residues, a share
$1/\gcd(n,M)$ that equals $1/n$ once $n\mid M$ and is never
smaller, so a fixed finite obstruction with moduli
$n_1,\dots,n_k$ touches every cylinder for small $M$ but, for
$M$ a multiple of their least common multiple, at most a share
$\sum_j1/n_j$ of them, which is below $\varepsilon$ when the
moduli are large though never $0$.  What the relaxed hypothesis
still demands, and what remains unproved in general, is that
the unresolved share can be driven to $0$ at all.
\end{remark}

\begin{remark}
Induced systems may repeat a modulus, whereas the systems of the
introduction have strictly increasing ones.  This costs nothing:
the proof of Theorem~\ref{thm:bounded} uses only the cylinder
selection and Lemma~\ref{lem:tail}, neither of which refers to
distinctness.

Theorem~\ref{thm:cyl} is not vacuous, and covers a system the earlier
results do not.  Take the doubled construction of Section~\ref{sec:routes}: its
misaligned parts are the small primes $q$ and so are unbounded,
and for any $M$ the classes excluded from
Corollary~\ref{cor:mixed} are those with $q\nmid M$, whose drift
$\sum_{q>Q}s/q$ diverges by Mertens -- so neither
Theorem~\ref{thm:bounded} nor Corollary~\ref{cor:mixed}
applies.  But at $M=2$ the odd cylinder carries no classes at
all, and the even cylinder's induced system has profinite
survivor measure $0$ -- its blocks are independent and the
product of their factors tends to $0$ -- so
Theorem~\ref{thm:cyl} gives its logarithmic density, which is
$1/2$.
\end{remark}

One hope for verifying the hypothesis is that refining $M$
might \emph{starve} cylinders of drift until they die.  It
cannot: refinement is drift-invariant.

\begin{proposition}[Drift invariance of refinement]\label{prop:driftinv}
For every modulus $M$, the average over the $M$ cylinders of the
induced drift equals the drift of the original system:
\[
  \frac1M\sum_{c\bmod M}\ \sum_{i\ \text{meeting}\ c}
  \frac{\gcd(n_i,M)}{n_i}\ =\ \sum_i\frac1{n_i}.
\]
\end{proposition}

\begin{proof}
Write $g=\gcd(n_i,M)$.  The class $(n_i,a_i)$ meets the cylinder
$c$ exactly when $g\mid a_i-c$, which happens for $M/g$ of the
$M$ residues, and on such a cylinder its induced modulus is
$n_i/g$, contributing induced drift $g/n_i$.  Its total
contribution over all cylinders is therefore
$(M/g)\cdot(g/n_i)=M/n_i$, independent of $M$.  Summing over $i$
and dividing by $M$ gives the claim.
\end{proof}

\begin{remark}
So refinement conserves drift on average: it cannot make every
cylinder drift-poor, though divergent drift may concentrate on
a small proportion of cylinders, as the doubled construction of
Section~\ref{sec:routes} shows with $M=2$.  Whatever verifies
the hypothesis of Theorem~\ref{thm:cylapp} must come from the
arithmetic of the residues, not from thinning the drift.  This
also explains why iterating the decomposition does not obviously
terminate: the object being refined is invariant under the
refinement.
\end{remark}

A useful special case identifies the cylinders by a coprimality
criterion rather than by inspection.

\begin{corollary}\label{cor:dense}
Suppose there is a modulus $M$ such that, for every residue $c$
modulo $M$, the induced system at $c$ contains a pairwise-coprime
subfamily of divergent drift.  Then $\dlog(A)$ exists and equals
$0$.
\end{corollary}

\begin{proof}
Fix $c$.  By Theorem~\ref{thm:D} applied to the induced system
at $c$, its survivor set has density $0$; the same product
$\prod(1-1/n)\to0$ is its profinite survivor measure, which is
therefore $0$ as well.  So every cylinder satisfies the first
alternative of Theorem~\ref{thm:cyl}, giving that $\dlog(A)$
exists, and each contributes $0$.
\end{proof}

\begin{remark}
This settles the systems on which
Proposition~\ref{prop:spike} is vacuous.  Take the moduli to be
the even numbers, each carrying two classes, one with an even
and one with an odd residue, the residues otherwise in
\emph{general position} -- here and below, residues drawn
independently and uniformly at random, statements about which
are measurements of one such draw unless a theorem is cited.  No two moduli are
coprime, so Theorem~\ref{thm:D} does not apply to the system
itself; but a class $(2m,a)$ with $a$ even kills only even $x$,
and writing $x=2y$ its condition becomes $y\equiv b\pmod m$, so
on the even cylinder the induced moduli run over \emph{all}
integers $m$ and contain the primes.  The classes with $a$ odd do
the same on the odd cylinder.  Both cylinders then satisfy the
hypothesis and the survivors vanish, as measured.

The hypothesis also correctly excludes the aligned families,
which survive: for the even moduli at residue $0$ the odd
cylinder carries no classes at all, so it has no coprime
subfamily of divergent drift, and the corollary does not apply --
as it must not, since those survivors have density $1/2$.
\end{remark}

\section{Reduction theorems}\label{sec:reductions}

The results so far resolve a system cylinder by cylinder.  The
theorems of this section reduce a system to another one instead,
and are the paper's principal positive results.  A first
mechanism exhausts cylinders rather than resolving them.

\begin{theorem}[Exhaustion as reduction]\label{thm:exhaust}
Suppose the classes partition into groups $G_1,G_2,\dots$, where
group $G_j$ kills only integers in a residue class
$C_j=c_j+d_j\mathbb{Z}$ (its \emph{protector}), and suppose the induced system of
$G_j$ on $C_j$ has profinite survivor measure $0$ for every $j$.
Call $\{(d_j,c_j)\}$ the \emph{reduced system}, $A'$ its
survivor set and $\nu'$ its survivor measure.  If $\dlog(A')$
exists and equals $\nu'$, then $\dlog(A)$ exists and equals
$\nu'=\nu$.
\end{theorem}

\begin{proof}
Write $U'$ for the killed set of the reduced system, each
reduced class $(d_j,c_j)$ being given the threshold
$\min(d_j,\min_{i\in G_j}t_i)$, at most $d_j$ and, for all but
the finitely many groups containing an exceptional class, at
least $d_j/M'$, so that the reduced system is a threshold system
with the same constant.  Since group $j$ kills only inside $C_j$ and only
from its thresholds on, $U\subseteq U'$, so
$\overline{\delta}_{\log}(U)\le\overline{\delta}_{\log}(U')$.
Since the profinite unions of the two systems agree up to a
null set, $\nu'=\nu$, and by hypothesis
$\dlog(U')=1-\nu'=\lim_k\delta\bigl(\bigcup_{j\le k}C_j\bigr)$,
the last equality because the finite unions have measures
increasing to $1-\nu'$.  For
the lower bound fix $k$: by Corollary~\ref{cor:nuzero} applied
to the induced system on $C_j$, group $j$'s survivors inside
$C_j$ have logarithmic density $0$, so
$U\supseteq\bigcup_{j\le k}C_j$ up to a set of logarithmic
density $0$, giving
$\underline{\delta}_{\log}(U)\ge\delta\bigl(\bigcup_{j\le k}C_j\bigr)$
for every $k$.  Letting $k\to\infty$ matches the upper bound.
\end{proof}

\begin{remark}
The aligned case -- all $c_j\equiv c$ for one integer
$c\ge0$ -- is the reduced system of Theorem~\ref{thm:transl};
the sparse case $\sum1/d_j<\infty$ is that of
Theorem~\ref{thm:conv}.  The theorem may be iterated: the
reduced system may itself be exhausted by its own grouping,
reducing further.  A \emph{tower} of depth two is exhibited
below, evading the single-level form.
\end{remark}

\begin{remark}
This theorem is not a curiosity: it covers a system that
\emph{none} of the others reach, and which refutes the naive
reduction discussed in Section~\ref{sec:reduction}.  Take sparse primes $p_j$ with
$\sum1/p_j<\infty$, and let $G_j$ consist of the classes of
modulus $p_jm$, for $m\ge2$ coprime to $p_j$, with residue
$\equiv1\pmod{p_j}$ and $\equiv0\pmod m$.  Then $G_j$ kills only
inside $C_j=\{x\equiv1\ (p_j)\}$, and within it kills every
point divisible by some $m\ge2$ coprime to $p_j$; the survivors
in $C_j$ are divisible by no prime other than $p_j$, a set of
measure $\prod_{q\ne p_j}(1-1/q)=0$, so $G_j$ exhausts $C_j$.
The moduli have greatest common divisor $1$ and admit no finite
prime cover, since $p_j$ itself is a modulus of $G_j$; the
residues are misaligned; each group's drift diverges; and yet
$\nu=\prod_j(1-1/p_j)>0$.  The same holds, as a measurement,
when the residues modulo $m$ are instead taken in general
position: with $p_j\in\{11,197,439,691,977,1259\}$ the
survivor density in the window $[X/2,X]$, $X=4\cdot10^6$, is
$0.8995$, agreeing with the product to four decimals; in the
three cylinders measured the survivor share is $0.0000$, and in
their complement $1.0000$.  So positive survivor measure against
divergent misaligned drift does \emph{not} require a common
factor: many protectors, one per group, suffice.

Nor need the protecting cylinders be aligned.  Take three sparse
primes $q_i$, here $11,197,439$; for each, level-one classes $(q_im,c)$ with
$c\equiv1\pmod{q_i}$ and $c$ in general position modulo $m$ --
a misaligned family of $172$ cylinders -- and let each be
exhausted by its own group of level-zero classes.  The
level-zero system has gcd $1$, drift $2.26$, and $20{,}488$
classes; its exhausted cylinders are not aligned, so a single
application of Theorem~\ref{thm:exhaust} with the aligned
reduced case does not reach it.  Two applications do: the
level-zero system reduces to the level-one family, which reduces
in turn to $\{(q_i,1)\}$, of convergent drift.  The measured
survivor density is $0.9024$, agreeing with
$\prod(1-1/q_i)$ to four decimals, the six level-one cylinders measured
having survivor share $0.0000$.
\end{remark}

Exhaustion suggests a majorant that needs no partition into
groups.

\begin{theorem}[Saturation]\label{thm:sat}
Let the system be standard ($t_i=n_i$), and let $\mathcal{P}$ be the family of cylinders $C$ with
$\mu(\Ahat\cap C)=0$ that are maximal under inclusion.  Then every
class lies inside a member of $\mathcal{P}$; the union of
$\mathcal{P}$ is the union of the balls up to a null set, so the
system $\mathcal{P}$ has survivor measure $\nu$; $\mathcal{P}$ is
\emph{saturated}, in that no proper coarsening of a member has
null survivor measure while every cylinder of null survivor
measure lies inside a member; and $\mathcal{P}$ admits no
proper null coarsening, so that the reduction applied to
$\mathcal{P}$ returns $\mathcal{P}$.  If the survivor set of $\mathcal{P}$ has
logarithmic density equal to $\nu$, then so does $A$.
\end{theorem}

\begin{proof}
The cylinders containing a given cylinder $c+d\Zhat$ are its
coarsenings $c+d'\Zhat$ with $d'\mid d$, finitely many; so among
the coarsenings of a null cylinder that are themselves null there
is a maximal one, and it is maximal under inclusion among all
null cylinders, hence a member of $\mathcal{P}$.  Applied to a
class $B_i$, which is null since $\mu(\Ahat\cap B_i)=0$, this puts
every class inside a member; applied to an arbitrary null
cylinder it gives the second half of saturation, and the first
half is the maximality of the members.  Every member of
$\mathcal{P}$ meets $\Ahat$ in a null set and so lies in the union
of the balls up to a null set, which gives the union and the
measure.  The profinite survivor set of $\mathcal{P}$ is $\Ahat$ up
to a null set, and a null cylinder for one set is a null cylinder
for the other, since
$|\mu(C\cap\Ahat)-\mu(C\cap\Ahat')|\le\mu(\Ahat\,\triangle\,\Ahat')$;
so the two sets have the same maximal null cylinders and the
reduction is idempotent.  For the density, a class inside $P\in\mathcal{P}$
kills only integers at least its modulus, which is a multiple of
$d_P$, while $P$ kills every integer at least $d_P$ in its
residue class; so the kills of $\mathcal{P}$ contain those of the
original, $A\supseteq A_{\mathcal{P}}$, and
$\underline{\delta}_{\log}(A)\ge\dlog(A_{\mathcal{P}})=\nu$, with
Theorem~\ref{thm:upper} for the other inequality.
\end{proof}

\begin{remark}
Unlike Theorem~\ref{thm:exhaust}, no partition of the classes
into groups is involved -- the members of $\mathcal{P}$ may meet
without nesting -- and the density statement is one-sided: the
saturated system majorises the original, and a density for it
passes down.  The consequence is that in proving $\dlog(A)=\nu$
one may assume the system saturated: its classes are the maximal
cylinders meeting the survivor set in a null set -- ``holes'' up
to measure zero, not necessarily disjoint from it -- no proper
coarsening is null, and any cylinder that dies is already inside
a class.  The
converse direction, that a density for $A$ forces one for
$A_{\mathcal{P}}$, we have not proved and do not need.  On
finite systems modulo $1260$ the reduction was computed exactly:
a many-protector system of $1{,}062$ classes collapses to its two
protectors, a coprime family of four primes is fixed, and on six
random systems every class lies in a member, the union is
reproduced, the reduction is idempotent and its output saturated.
\end{remark}

In the saturated form the classes are the maximal null cylinders of $\Ahat$, and
Lemma~\ref{lem:escape}, read backwards, says what positive
survivor measure forces of them.

\begin{corollary}[Anchoring]\label{cor:anchor}
Let the system be saturated with $\nu>0$, and for each $P$ let
$\mathcal{H}_P$ be the classes whose moduli are coprime to every
prime at most $P$, and let $\rho_P$ be the limsup of the
second-moment ratios of $\mathcal{H}_P$ along the exhaustion by
increasing modulus.  Then either $\sum_{\mathcal{H}_P}\mu(C)<\infty$
for some $P$, in which case every class outside
$\mathcal{H}_P$ lies in one of the finitely many cylinders
$c+p\mathbb{Z}$ with $p\le P$; or $\sup_P\rho_P=\infty$.
\end{corollary}

\begin{proof}
A class in $\mathcal{H}_P$ is determined by the coordinates at
primes above $P$, and its points are all killed.  If the sums
diverged for every $P$ and the $\rho_P$ were bounded by one
constant, Lemma~\ref{lem:escape} would put almost every point in
infinitely many classes, so $\nu=0$.  The alternative in the
first case is the definition of $\mathcal{H}_P$: a modulus not
coprime to the primes up to $P$ is divisible by one of them.
\end{proof}

\begin{remark}
So a saturated system with survivors either splits, at some
level $P$, into escapees of summable measure -- the summability
hypothesis of Theorem~\ref{thm:confined}, though not its
covering hypothesis on finite unions -- and classes anchored to
a bounded family of prime cylinders, or its escapees are
concentrated in a sense that strengthens without bound as $P$
grows.  Writing $N(x)$ for the
number of classes containing $x$, the second alternative says
$\int N^2\gg(\int N)^2$ along the escapees: a set of small
measure lies in an enormous number of holes.  A counterexample
must therefore either exhibit such concentration, or have its
anchored classes -- an infinite system confined to finitely many
prime cylinders, possibly of divergent drift -- fail the
covering hypothesis of Theorem~\ref{thm:confined} there.  We
have not been able to convert either into one of the
organisations covered above.
\end{remark}

Exhaustion asks the groups to die.  They need not, if the
cylinders confining them are sparse.

\begin{theorem}[Summable protectors]\label{thm:confined}
Suppose the classes partition into groups $G_1,G_2,\dots$, where
group $G_j$ kills only integers in a residue class
$C_j=c_j+d_j\mathbb{Z}$, that $\sum_j1/d_j<\infty$, and that for
every $J$ the survivor set $A_J$ of $G_1\cup\dots\cup G_J$ has
logarithmic density equal to its profinite survivor measure
$\nu_J$.  Then $\dlog(A)$ exists and equals $\nu$.
\end{theorem}

\begin{proof}
The profinite unions of the finite systems increase to that of
the full system, so $\nu_J\downarrow\nu$, and by hypothesis
$\dlog(A_J)=\nu_J$.  Since $A\subseteq A_J$,
$\overline{\delta}_{\log}(A)\le\nu_J$ for every $J$, hence
$\le\nu$.  Taking $J$ past the finitely many groups containing an
exceptional class, the groups beyond $J$ kill only inside their
cylinders, and only integers at least $d_j/M'$ by the fourth
threshold fact ($M'=1$ for a standard system); so
$A\supseteq A_J\setminus\bigcup_{j>J}\bigl(C_j\cap[d_j/M',\infty)\bigr)$.
A residue class modulo $d$ contributes at most
$(\log X+\log M'+M')/d$ to the harmonic sum over $[d/M',X]$, so the union has
upper logarithmic density at most $\sum_{j>J}1/d_j$, and
$\underline{\delta}_{\log}(A)\ge\nu_J-\sum_{j>J}1/d_j$, which
tends to $\nu$.
\end{proof}

\begin{remark}
Theorem~\ref{thm:conv} is the case of singleton groups, each
confined to its own ball, and Theorem~\ref{thm:exhaust} with
summable protectors is the case of dying groups.  The point is
that the groups need not die: Proposition~\ref{prop:irred2}
below has divergent drift, positive survivor measure and no
dying cylinder whatever, and this theorem covers it.
\end{remark}

Both reductions are instances of one structure, and that
structure can be pushed to its limit.  A \emph{tree of
cylinders} is a family $T$ of residue classes $c+d\mathbb{Z}$,
$d\ge1$, containing $\mathbb{Z}$, any two of which are nested or
disjoint, each member having finitely many maximal proper
sub-members, its \emph{children}; a member's \emph{gap} is what
its children leave uncovered, and its \emph{depth} is the number
of its proper ancestors.  A system is \emph{coprime-layered} on
$T$ if
\begin{enumerate}
\item[(T1)] $T$ is a tree of cylinders as just defined;
\item[(T2)] every class is $(d_Cq,a)$ with $C=c+d_C\mathbb{Z}$ in
$T$, $q$ prime and $a\equiv c\pmod{d_C}$: the class is
\emph{attached} to $C$ and lies inside it;
\item[(T3)] no class is contained in a class attached to a
proper ancestor of its node;
\item[(T4)] no attached prime divides the modulus of any node.
\end{enumerate}
Write $S_C=\sum1/q$ over the classes attached to $C$.  The
nodes containing a point $x\in\Zhat$ form a chain, its
\emph{branch}, and $g(x)=\sum_CS_C$ over that chain is the
\emph{branch drift}; $g_D$ denotes the sum over the nodes of
depth at most $D$.  Nothing here restricts the drift, the
survivor measure, the alignment, or the depth.

\begin{lemma}[Dying branches]\label{lem:dying}
In a coprime-layered system, $\mu(\{g=\infty\}\cap
\{\text{survivors}\})=0$.
\end{lemma}

\begin{proof}
Fix a node $C$ of depth $D$ and a point $x\in C$ with index $t$,
$x=c+d_Ct$.  A class attached to an ancestor $C'$ of $C$, or to
$C$ itself, restricts on $C$ to a single residue class of $t$
modulo its prime, because by (T4) the prime is coprime to
$d_C/d_{C'}$; by (T3) two classes on the chain with the same
prime restrict to different residues, and so do two classes at
the same node with the same prime, being distinct balls with the
same residue modulo $d_C$.  The index $t$ is
uniformly distributed modulo any product of these primes, so the
points of $C$ surviving the classes on the chain to $C$ have
measure $\mu(C)\prod_q(1-k_q/q)\le\mu(C)e^{-g_D}$, with $k_q$
the number of classes of prime $q$ on the chain.  The set
$\{g_D>N\}$ is a union of nodes of depth at most $D$ and of
gaps, on each of which the same bound holds with the chain to
that cell, so $\mu(\{g_D>N\}\cap\{\text{survivors}\})\le e^{-N}$.
Since $\{g=\infty\}\subseteq\bigcup_D\{g_D>N\}$ for every $N$,
and the sets on the right increase with $D$, the survivors in
$\{g=\infty\}$ have measure at most $e^{-N}$ for every $N$.
\end{proof}

\begin{lemma}[Light and heavy]\label{lem:light}
In a coprime-layered system, for every $N>0$,
$\underline{\delta}_{\log}(A)\ge\nu-\mu\{g>N\}$.
\end{lemma}

\begin{proof}
Write $g_C$ for the branch drift summed to $C$ inclusive.  Call a
class \emph{light} if $g_C\le N$ at its node, and call a node a
\emph{first crossing} if $g_C>N$ while every proper ancestor has
$g\le N$.  First crossings are pairwise disjoint, being
non-nested members of $T$, and their union is exactly
$\{g>N\}$.  The light classes have total drift
$\sum_C\mu(C)S_C$ over nodes with $g_C\le N$, which is
$\int\sum_{C\ni x,\ g_C\le N}S_C\,d\mu(x)\le N$, so by
Theorem~\ref{thm:conv} their survivor set $A'$ has logarithmic
density $\nu'\ge\nu$.  Every other class is attached to a node
at or below a first crossing $C_j$, lies inside $C_j$, and kills
only integers at least $d_j/M'$, by the fourth threshold fact.
Hence $A\supseteq A'\setminus\bigcup_j(C_j\cap[d_j/M',\infty))$,
and as in Theorem~\ref{thm:confined}
$\underline{\delta}_{\log}(A)\ge\nu'-\sum_j1/d_j
=\nu'-\mu\{g>N\}\ge\nu-\mu\{g>N\}$.
\end{proof}

\begin{theorem}[Coprime-layered trees]\label{thm:tree}
For every coprime-layered system, $\dlog(A)$ exists and equals
$\nu$.
\end{theorem}

\begin{proof}
The upper bound is Theorem~\ref{thm:upper}.  For the lower bound
fix $\varepsilon>0$ and $0<\eta<\tfrac12$.  Let $\mathcal{F}_D$ be the finite
partition of $\Zhat$ into the nodes of depth $D$ and the
residue classes composing the gaps of nodes of smaller depth,
its \emph{cells}; each cell is a single residue class, and the
partitions refine as $D$ grows, since a gap once formed is never
subdivided.  The dying set $E=\{g=\infty\}$ is determined
by the branch, hence measurable for
$\mathcal{F}_\infty=\sigma(\bigcup_D\mathcal{F}_D)$, so by
L\'evy's upward theorem~\cite{Williams} the conditional measure
of $E$ in the cell containing $x$ tends to $1_E(x)$ for almost
every $x$.  Choose $D$ so large that the cells in which the
conditional measure of $E$ lies in $(\eta,1-\eta)$ -- the
\emph{mixed} cells -- have total measure at most $\varepsilon$;
call the remaining cells \emph{dying} or \emph{alive} according
as that measure is at least $1-\eta$ or at most $\eta$.

Let $C$ be an alive cell, a residue class $c+d_C\mathbb{Z}$;
either $C$ is a node of depth $D$, or $C$ is one of the residue
classes composing the gap of a node $C_0$, in which case its
modulus is a multiple of $d_{C_0}$ built from node moduli.  Write
$x=c+d_Ct$.  The classes that meet $C$ are those attached to the
nodes containing $C$ (its node and that node's ancestors) and,
when $C$ is a node, those attached to its descendants.  By (T4)
every attached prime $q$ is coprime to $d_C$, so a class
$(d_{C'}q,a)$ attached to a node $C'\supseteq C$ becomes in $t$
the class of modulus $q$, attached to the root, and a class
attached to a descendant $C''\subset C$ becomes one of modulus
$(d_{C''}/d_C)\,q$, attached to the node $C''$ of the subtree
rooted at $C$.  This is a coprime-layered system on that subtree
(a single root when $C$ is a gap class): (T1) and (T4) are
inherited; (T2) holds with the induced node moduli; and (T3)
holds because two root classes with the same prime $q$ and the
same residue modulo $q$ would come from classes agreeing modulo
$q$ and, both containing $C$, modulo their node moduli, so that
one would contain the other, which (T3) for the original system
forbids.  Its branch drifts are those of the original along the
branches through $C$, its survivor set is the index set of
$A\cap C$, and its survivor measure is
$\nu_C=\mu(C\cap\{\text{survivors}\})/\mu(C)$.  By
Lemma~\ref{lem:restrict} it is a threshold system, and
Lemma~\ref{lem:light} applies to it by the fourth threshold fact.
Lemma~\ref{lem:light} gives
$\underline{\delta}_{\log}(A\cap C)\ge\mu(C)(\nu_C-\mu\{g>N\mid C\})$
for every $N$, and letting $N\to\infty$,
$\underline{\delta}_{\log}(A\cap C)\ge\mu(C)\nu_C-\eta\mu(C)$.
On a dying cell, Lemma~\ref{lem:dying} gives
$\mu(C\cap\{\text{survivors}\})\le\eta\mu(C)$.  Lower
logarithmic density is superadditive over disjoint sets, so
summing over the alive cells,
\[
\underline{\delta}_{\log}(A)\ \ge\
\sum_{\text{alive}}\mu(C\cap\{\text{survivors}\})-\eta
\ \ge\ \nu-\eta-\varepsilon-\eta,
\]
the survivors in dying cells having measure at most $\eta$ and
those in mixed cells at most $\varepsilon$.  Let $\eta$ and
$\varepsilon$ tend to $0$.
\end{proof}

\begin{remark}
The theorem needs neither dying groups nor summable protectors.
It covers, for instance, a tree whose deep cylinders retain
measure bounded away from $0$, whose nodes carry, in a
proportion $1/k^2$ at depth $k$, finite drifts so large that the
total drift diverges, and along almost all of whose branches the
drift nonetheless converges, so that a positive measure of
survivors lies at infinite depth: the case that
Theorems~\ref{thm:exhaust} and~\ref{thm:confined} both miss.
Its reach is nevertheless limited, and in a way that is easy to
misjudge: being coprime-layered is a strong hypothesis on the
moduli, not a normalisation.

\begin{proposition}[Chain constraint]\label{prop:chain}
In a coprime-layered system, any two nodes carrying classes have
non-coprime moduli, unless one of the moduli is $1$.
\end{proposition}

\begin{proof}
If $d,d'>1$ are coprime then, by the Chinese remainder theorem,
the classes $c+d\mathbb{Z}$ and $c'+d'\mathbb{Z}$ meet whatever
$c,c'$ are, so the two nodes are not disjoint and must be
nested; but nesting forces $d\mid d'$ or $d'\mid d$, which with
$\gcd(d,d')=1$ makes one of them $1$.
\end{proof}

So the classes of a coprime-layered system hang off a family of
cylinders any two of which share a factor: the node moduli
$2,10,100,1700$ of the explicit tree described in
Appendix~\ref{sec:verification}, not the moduli $7,37,73,\dots$
of a family of protectors.
Proposition~\ref{prop:irred2} is therefore \emph{not}
coprime-layered -- its natural nodes are the distinct primes
$p_j$, pairwise coprime -- and it is covered by
Theorem~\ref{thm:confined} instead.  Nor are the
semiprimes, whose nodes would be the primes $p$; they are
covered by Theorem~\ref{thm:quasi}.  Checked by exact search
over the legal node assignments on truncations -- for the tree,
its depth-two slice -- that explicit tree and the multiples of
$6$ are coprime-layered, while those two systems are not.
\end{remark}

\begin{remark}\label{rem:circular}
Lemma~\ref{lem:dying} is the only place where primality of the
attached moduli is used, which invites the thought that allowing
a node to carry a layer of composite moduli would extend the
theorem.  It would not: the lemma's analogue for composite layers
is false.  Proposition~\ref{prop:irred2} below is a system on the
single cylinder $\Zhat$, with composite moduli, divergent
drift and no proper reduction, whose survivor measure is
positive; read as a layer it has $g=\infty$ everywhere and yet
does not kill almost every point.  So without primality the
proof of Theorem~\ref{thm:tree} collapses at its first step, and
an extension would have to establish, for each layer, the
reverse inequality itself.  The tree theorem is exactly as
strong as the independence its coprimality provides.
\end{remark}

The chain constraint says which nested families the tree
theorem reaches.  At the opposite extreme, protectors with
pairwise coprime moduli, the difficulty disappears for a
different reason: positive survivor measure itself produces the
confinement.

\begin{theorem}[Coprime protectors]\label{thm:coprot}
Suppose the classes partition into groups $G_j$, where $G_j$
consists of classes $(p_jq,a)$ with $a\equiv c_j\pmod{p_j}$ and
$q$ prime, the moduli $p_j$ are pairwise coprime, the attached
primes $q$ are distinct across all classes, and none of them
divides any $p_i$.  Then $\dlog(A)$ exists and equals $\nu$,
whatever $\sum_j1/p_j$ may be.
\end{theorem}

\begin{proof}
Write $S_j=\sum1/q$ over the primes attached in $G_j$ and
$\kappa_j=1-\prod(1-1/q)$, the measure of the kills of $G_j$
inside $C_j=c_j+p_j\Zhat$ relative to $\mu(C_j)$.  Whether $x$
lies in $C_j$ depends on $x$ modulo $p_j$; whether $x\in C_j$ is
killed by $G_j$ depends on $x$ modulo the primes attached in
$G_j$.  By hypothesis these coordinates are distinct across $j$,
and the coordinates of $\Zhat$ are independent, so
\[
\nu=\prod_j\Bigl(1-\frac{\kappa_j}{p_j}\Bigr).
\]
If $\sum_j\kappa_j/p_j=\infty$ then $\nu=0$ and
Corollary~\ref{cor:nuzero} finishes.  Otherwise, since
$\kappa_j\ge1-e^{-S_j}\ge(1-e^{-1})\min(S_j,1)$, we get
$\sum_j\min(S_j,1)/p_j<\infty$: the groups with $S_j\ge1$ lie on
cylinders with $\sum1/p_j<\infty$, and the remaining groups have
total drift $\sum S_j/p_j<\infty$.  A group with $S_j=\infty$ is
a coprime layer of divergent drift, which dies on $C_j$ by
Theorem~\ref{thm:D}; apply Theorem~\ref{thm:exhaust} with those
groups and every other class as a singleton group on its own
ball, where its induced system is the single class of modulus
$1$ and dies trivially; this reduces each dying group to the
class $(p_j,c_j)$ and leaves every other class unchanged.  The reduced system is covered
by Theorem~\ref{thm:confined}, with the groups: each reduced
class and each group with $1\le S_j<\infty$, confined to its
$C_j$, the moduli having $\sum1/p_j<\infty$; and each remaining
class as a singleton confined to its own ball, these measures
summing to $\sum_{S_j<1}S_j/p_j<\infty$.  A finite union of
these groups has finite drift, so Theorem~\ref{thm:conv} gives
its density as its survivor measure, the hypothesis of
Theorem~\ref{thm:confined}.
\end{proof}

\begin{remark}
This is the shape the checklist of Section~\ref{sec:reduction}
isolates -- protectors
too dense to be summable, kills too sparse to be quasi-independent
-- and the theorem says it cannot occur with coprime protectors:
if the kills are sparse enough for survivors, the heavy
protectors are already summable and the light classes already
convergent.  On an instance with $61$ protectors, two of them
carrying layers of drift $0.40$ and the rest layers of drift
$1/(40\log^2p)$, all attached primes distinct, the product gives
$\nu=0.94479$ against a Monte Carlo value of $0.9413$.  Distinct
primes make a single layer of drift $1$ need primes beyond
$10^9$, so the instance checks the formula and the split, not
divergence; for that, a layer of drift $1.5$ on every odd prime
has $\kappa\ge1-e^{-1.5}$, and the exact partial products
$\prod_{p\le N}(1-\kappa/p)$ fall through $0.2492$ at $N=10^{3}$, $0.1997$ at $N=10^{4}$, $0.1680$ at $N=10^{5}$, $0.1459$ at $N=10^{6}$, to $0$ like
$(\log N)^{-0.78}$.  What the theorem does not reach is
protectors that share factors without nesting -- pairwise
products $p_ip_j$, say -- where membership is neither
independent, as here, nor a chain, as in Theorem~\ref{thm:tree}.
\end{remark}

That last case is the remaining one, and it is settled by
Lemma~\ref{lem:escape}.

\begin{theorem}[Overlapping protectors]\label{thm:overlap}
Let $\mathcal{S}$ be a set of primes and, for every pair
$p\ne p'$ in $\mathcal{S}$, let $G_{pp'}$ be a group of classes
killing only inside a cylinder $C_{pp'}$ of modulus $pp'$, whose
kills there have relative measure $\kappa_{pp'}$.  Assume that
each group is a coprime layer of primes outside $\mathcal{S}$,
distinct across groups, and that $\kappa_{pp'}\ge\kappa_0>0$.
If $\sum_{p\in\mathcal{S}}1/p$ diverges and every pair occurs,
then $\nu=0$ and $\dlog(A)=0$.  If it converges, then
$\sum1/(pp')<\infty$ and $\dlog(A)=\nu$ by exhaustion and
confinement.  Either way $\dlog(A)$ exists and equals $\nu$.
\end{theorem}

\begin{proof}
In the convergent case $\sum_{p<p'}1/(pp')\le\frac12(\sum1/p)^2$
is finite.  A group whose layer has divergent drift dies on its
cylinder by Theorem~\ref{thm:D}; apply Theorem~\ref{thm:exhaust}
to those groups, with every other class a singleton, reducing
each to the class of its cylinder.  Then apply
Theorem~\ref{thm:confined} to the reduced system with the groups
confined to the cylinders $C_{pp'}$, of summable measure: a
finite union of them has finite drift, so Theorem~\ref{thm:conv}
covers it.
Suppose then that $A:=\sum_{p\in\mathcal{S}}1/p=\infty$, and let
$\Kh_{pp'}$ be the kills of $G_{pp'}$, so
$\mu(\Kh_{pp'})=\kappa_{pp'}/(pp')$.  Delete from $\mathcal{S}$ the primes at most $P$ and both
endpoints of each of the finitely many groups whose attached
primes include one at most $P$ -- finitely many because each
prime is attached to at most one group.  Finitely many primes
are removed, so the reciprocal sum of the remaining primes still
diverges, every pair of remaining primes still carries its
group, and the remaining family is determined by the coordinates
above $P$; enumerate it by the pairs among the first $N$
remaining primes as $N$ grows.  For
the second moment, two pairs sharing no prime give
independent events, since the coordinates determining them --
the two protector primes and the attached primes of the group,
which lie outside $\mathcal{S}$ and are distinct across groups
-- are disjoint; and two sharing the prime $p$ give
$\mu(\Kh_{pp'}\cap \Kh_{pp''})\le\mu(C_{pp'}\cap C_{pp''})
=1/(pp'p'')$.  Writing $A_N,B_N$ for $\sum1/p$ and $\sum1/p^2$
over the first $N$ primes of $\mathcal{S}$ above $P$, and
$M_N=(A_N^2-B_N)/2$ for the total measure of the cylinders,
\[
\sum_{e,f}\mu(C_e\cap C_f)=M_N+\sum_p\frac1p\Bigl[(A_N-1/p)^2-(B_N-1/p^2)\Bigr]
+\Bigl(M_N^2-\Sigma'\Bigr),
\]
where $e,f$ run over ordered pairs of edges of the complete
graph on those $N$ primes, $\mu(C_e)=1/pp'$ for $e=\{p,p'\}$,
and the three terms are as follows.  The diagonal $e=f$
contributes $\sum_e\mu(C_e)=M_N$.  Two distinct edges
$e=\{p,p'\}$, $f=\{p,p''\}$ sharing exactly the prime $p$ meet
in the cylinder of modulus $pp'p''$, and summing $1/pp'p''$ over
ordered pairs $p'\ne p''$ of primes other than $p$ gives
$p^{-1}[(A_N-1/p)^2-(B_N-1/p^2)]$; summing over $p$ gives the
middle term.  Disjoint edges have independent cylinders, so their
contribution is $\sum_{e\cap f=\varnothing}\mu(C_e)\mu(C_f)$,
which is $M_N^2$ less the products over non-disjoint ordered
pairs, $\Sigma'=\sum_e\mu(C_e)^2+\sum_pp^{-2}[(A_N-1/p)^2-(B_N-1/p^2)]$.
Since $\Kh_e\subseteq C_e$, the numerator for the kills is at
most this sum; since $\mu(\Kh_e)\ge\kappa_0\mu(C_e)$, the
denominator is at least $\kappa_0^2M_N^2$.  The middle sum is at
most $A_N^3$, $\Sigma'\le B_N^2+A_N^2B_N$, and $M_N\sim A_N^2/2$,
so the displayed sum is $M_N^2(1+O(1/A_N))$, and the
second-moment ratio of the kills is at most
$\kappa_0^{-2}(1+O(1/A_N))$.  So
Lemma~\ref{lem:escape} applies with $\rho=2\kappa_0^{-2}$, almost
every point lies in infinitely many $\Kh_e$, and $\nu=0$;
Corollary~\ref{cor:nuzero} finishes.
\end{proof}

\begin{remark}
The ratio's approach to $1$ is genuinely slow, since
$A_N\asymp\log\log N$: over all primes up to
$10^3,10^4,10^5,10^6,10^7$ the exact ratios are 3.72, 3.29, 3.03, 2.87, 2.74, while the
normalised quantity $(R-1)A/4$ falls through 1.1553, 1.1345, 1.1217, 1.1129, 1.1065 toward the
predicted $1$.  Over the same range the measure of the points
lying in no protector at all falls through 0.4811, 0.3965, 0.3392, 0.2975, 0.2658, and at $N=10^3$
agrees with a Monte Carlo value of $0.4767$.  The sparse
instance $p_i$ the least prime at least $i^2$ has, at $49$
primes, total protector measure $0.1191$ and protector-free
measure $0.9070$ against a Monte Carlo value of $0.9044$: there the sums
converge and confinement, not this theorem, is what applies.
\end{remark}

\section{The reduction}\label{sec:reduction}

Recall $U=\bigcup_i K_i\subseteq\N$.  For $k\ge1$ let
$U_k=\bigcup_{i\le k}K_i$ and
\[
  F(X)=\frac{1}{\log X}\sum_{x\in U,\ x\le X}\frac1x,
  \qquad
  F_k(X)=\frac{1}{\log X}\sum_{x\in U_k,\ x\le X}\frac1x .
\]

\begin{theorem}\label{thm:reduction}
If
\[
  \sup_{X\ge2}\bigl(F(X)-F_k(X)\bigr)\longrightarrow 0
  \qquad (k\to\infty),
\]
then $\dlog(A)$ exists and equals $\nu$.
\end{theorem}

\begin{proof}
Each $U_k$ is a finite union of truncated arithmetic
progressions, hence periodic apart from finitely many integers,
so $F_k(X)\to 1-\nu_k$ as $X\to\infty$.  Let $\varepsilon>0$ and
choose $k$ with $\sup_{X\ge2}\bigl(F(X)-F_k(X)\bigr)<\varepsilon$.
Since $U_k\subseteq U$ we have, for every $X$,
\[
  F_k(X)\ \le\ F(X)\ \le\ F_k(X)+\varepsilon .
\]
Letting $X\to\infty$ along the upper and lower limits gives
\[
  1-\nu_k\ \le\ \liminf_X F(X)
  \ \le\ \limsup_X F(X)\ \le\ 1-\nu_k+\varepsilon,
\]
so $\limsup_X F-\liminf_X F\le\varepsilon$.  As $\varepsilon$
was arbitrary, $\lim_X F(X)$ exists; and since it lies within
$\varepsilon$ of $1-\nu_k$ for every such $k$, while
$\nu_k\downarrow\nu$, that limit is $1-\nu$.  Hence
$\dlog(A)=1-\lim_X F(X)=\nu$.
\end{proof}

It is worth recording an equivalent form in terms of window
densities, since it moves the difficulty from an infinite tail
to a set of scales.  For $X\ge2$ let
\[
  w(X)=\frac{\bigl|U\cap[X/2,X]\bigr|}{\bigl|\N\cap[X/2,X]\bigr|}
\]
be the killed density in the dyadic window at $X$.

\begin{lemma}[Window lower bound]\label{lem:window}
$\liminf_{X\to\infty}w(X)\ge1-\nu$.
\end{lemma}

\begin{proof}
Fix $Y$.  The classes of modulus at most $Y$ kill a periodic set
of density $1-\nu_Y$ whose period $L_Y$ is the lcm of those
moduli.  A window of length $X/2$ contains
$\lfloor X/(2L_Y)\rfloor$ full periods and a remainder shorter
than $L_Y$, so the density of that periodic set inside the
window is $1-\nu_Y+O(L_Y/X)$.  Since $U$ contains it,
$w(X)\ge1-\nu_Y-O(L_Y/X)$.  Let $X\to\infty$ with $Y$ fixed to
get $\liminf_X w\ge1-\nu_Y$, then let $Y\to\infty$ and use
$\nu_Y\downarrow\nu$.
\end{proof}

The reverse fails: classes of modulus close to $X$ kill one
element of the window apiece, so $w$ can reach $1$ while
$1-\nu$ is far below (Section~\ref{sec:routes}).  Thus $w$ is pinned from below
and free above, and the whole question is how often it escapes.

For a set $E$ of real scales $X\ge1$ define its logarithmic
density as the limit, when it exists, of
$\frac{1}{\log Y}\int_{E\cap[1,Y]}dX/X$, with upper and lower
versions as usual; $w(X)$ is defined for every real $X\ge2$ by
the window $[X/2,X]\cap\N$.

\begin{theorem}[Excursion form]\label{thm:excursion}
Writing $u=\log X$ and $w(v)$ for $w(e^v)$, one has
$F(e^u)=\frac{1}{u}\int_{v=\log 2}^{v=u} w(v)\,dv+O(1/u)$, so $\dlog(A)$ exists
and equals $\nu$ if and only if, for every $\delta>0$, the set
\[
  E_\delta=\{X:\ w(X)>1-\nu+\delta\}
\]
has logarithmic density $0$ in the sense just defined.
\end{theorem}

\begin{proof}
With $S(X)=\sum_{x\in U,\ x\le X}1/x$ and $Y=e^u$, integrate
$w$ against $dX/X$ over $[2,Y]$.  A killed integer $x\le Y/2$
lies in the window $[X/2,X]$ exactly for $X\in[x,2x]$, and the
window then holds $X/2+O(1)$ integers, so its weight is
$\int_x^{2x}\frac{dX}{X}\cdot\frac{1}{X/2+O(1)}
=\frac1x+O(x^{-2})$; a killed integer in $(Y/2,Y]$ receives at
most $\int_x^{Y}\frac{2\,dX}{X^2}\le\frac1x$, and there are at
most $Y/2$ of them, each weighted at most $2/Y$, so they
contribute $O(1)$.  Summing, $\int_{\log2}^u w(v)\,dv
=S(Y)+O(1)$, the errors $O(x^{-2})$ having a convergent sum.
(The relation is not an exact identity: for $U=\N$ one has
$w\equiv1$ while $F=1+\gamma/\log X+o(1/\log X)$.)  Hence
$F(e^u)=\frac{1}{u}\int_{v=\log 2}^{v=u} w(v)\,dv+O(1/u)$, whose error vanishes, so $F$
converges precisely when $w$ has a logarithmic mean.  By
Lemma~\ref{lem:window} the function $w-(1-\nu)$ is
non-negative in the limit, so its average vanishes if and only
if it is small off a set of vanishing logarithmic density, which
is the stated condition.  Finally $\dlog(A)=1-\lim F$.
\end{proof}

One part of the excess admits a quantitative bound.  Write
$D(Y)=\sum_{n_i\le Y}1/n_i$ and
$h=\limsup_{Y\to\infty}D(Y)/\log Y$, the harmonic density of the
modulus set; note that $h\le1$ when the moduli are distinct.

\begin{proposition}[Spike bound]\label{prop:spike}
Let $\sigma(X)$ denote the density in the window $[X/2,X]$ of
the union of the kills of the classes of modulus in $(X/2,X]$,
which bounds the contribution of those classes to $w(X)$ however
the overlaps are apportioned.  Then for every $\delta>0$ the set
$\{X:\ \sigma(X)>\delta\}$ has logarithmic density at most
$(2\log2)\,h/\delta$.  In particular it has logarithmic density
$0$ whenever $h=0$.
\end{proposition}

\begin{proof}
A class of modulus $n\in(X/2,X]$ kills at most one element of
$[X/2,X]$, since its elements are at least $n$ and consecutive
ones differ by $n>X/2$.  Hence, with $N$ the number of such
classes and $X\ge4$, $\sigma(X)\le N/(X/2-1)\le(2N/X)(1+4/X)$,
the window holding at least $X/2-1$ integers.  Each of these
classes contributes at least $1/X$ to the window drift
$d(X)=\sum_{X/2<n_i\le X}1/n_i$, giving $N/X\le d(X)$ and
therefore $\sigma(X)\le2(1+4/X)\,d(X)$.

Since $n_i\in(X/2,X]$ exactly when $X\in[n_i,2n_i)$,
\[
\int_1^Y d(X)\,\frac{dX}{X}
=\sum_{n_i<Y}\frac1{n_i}\int_{n_i}^{\min(2n_i,Y)}\frac{dX}{X}
\le(\log2)\,D(Y),
\]
so for $\delta'<\delta$ Markov's inequality for the measure
$dX/X$ on $[1,Y]$ bounds the $dX/X$-measure of
$\{X\le Y:\ 2d(X)>\delta'\}$ by $(2\log2)D(Y)/\delta'$, whose
ratio to $\log Y$ has limsup $(2\log2)h/\delta'$.  Since
$\sigma\le2(1+4/X)d$, the set $\{\sigma>\delta\}$ lies in
$\{2d>\delta'\}$ beyond a bounded $X$; let $\delta'\to\delta$.
\end{proof}

\begin{remark}\label{rem:intermediate}
The scope is exactly the spike term.  Classes of \emph{inter\-mediate}
modulus, say $T<n_i\le X/2$, are not controlled this way:
bounding their contribution by $\sum1/n_i$ over that range gives
about $h\log X$, and restricting to the survivors of the
small-modulus classes improves this only to $\nu_T\,h\log X$.
Both are useless.  What Proposition~\ref{prop:spike} does show
is that the mechanism of Section~\ref{sec:routes} -- fresh classes of modulus
comparable to $X$ -- can produce excursions on a set of positive
logarithmic density only when the moduli have positive harmonic
density.

The natural alternative also fails, and for a classical reason
worth recording.  One might hope that the window $[X/2,X]$
samples the union of the intermediate classes at its global
density.  It does not: that union is periodic with period the
lcm of the moduli involved, of size roughly $e^{X/2}$, vastly
longer than the window.  The discrepancy is not hypothetical.
Take the intermediate classes to be the primes $p\le X/2$ with
residue $0$.  Their survivors inside $[X/2,X]$ are exactly the
primes there, of density $\sim1/\log X$, whereas the global
survivor density of the same finite system is
$\prod_{p\le X/2}(1-1/p)\sim e^{-\gamma}/\log X$ by Mertens.  The
window therefore over-represents survivors by a factor tending
to $e^{\gamma}=1.781\ldots$; numerically the ratio runs
$1.7213$, $1.7278$, $1.7342$, $1.7388$ for
$X=2\cdot10^5,\,10^6,\,4\cdot10^6,\,1.6\cdot10^7$.  This is the
Buchstab phenomenon, and it shows that no argument based on the
window inheriting global densities can control the intermediate
range.  We regard this as the precise reason that range resists,
rather than as evidence about the answer.

One might hope that the discrepancy always runs in the harmless
direction.  In the example above the window retains \emph{more}
survivors than the global density predicts, so the intermediate
classes \emph{under}-kill, and under-killing cannot produce an
excursion.  If that were general, Proposition~\ref{prop:spike}
would control the whole excess.  It is not general.  Taking the
classes to be the primes in $[X^{a},X^{b}]$ at residue $0$, for
which coprimality makes the global survivor density the exact
product $\prod(1-1/p)$, the ratio of the window survivor density
to the global one is measured at $X=8\cdot10^6$ as
\[
  0.9949\ \ (a,b)=(0.15,0.35),\qquad
  0.9107\ \ (a,b)=(0.05,0.50),
\]
both below $1$, against $1.0174$ and $1.0549$ for
$(0.20,0.45)$ and $(0.25,0.50)$.  So the ratio crosses $1$ in
both directions, as the oscillation of Buchstab's function
suggests, and intermediate classes can genuinely over-kill.  The
spike bound alone therefore does not control the excess.
\end{remark}

The excess admits a clean expression in terms of the quantity
sieve theory studies.  Let $s(X)$ be the survivor density in the
window and $\nu_X$ the profinite survivor measure of the classes
of modulus at most $X$, and set $r(X)=s(X)/\nu_X$ -- the
window-to-global survivor ratio, the Buchstab ratio of
Remark~\ref{rem:intermediate}.

\begin{proposition}[Reduction to the Buchstab ratio]\label{prop:ratio}
Assume $\nu>0$.  For every $X$,
\[
  w(X)-(1-\nu)\ \le\ \nu\,\bigl(1-r(X)\bigr)^{+},
\]
so there is no excursion at all when $r(X)\ge1$.  Consequently
$E_\delta\subseteq\{X:\ r(X)<1-\delta/\nu\}$, and $\dlog(A)$
exists and equals $\nu$ as soon as $(1-r)^{+}\to0$ in
logarithmic mean; since $\limsup_Xr(X)\le1$ (shown below), this
is the same as $r\to1$ in logarithmic mean.
\end{proposition}

\begin{proof}
The excess is $w(X)-(1-\nu)=\nu-s(X)=\nu-r(X)\nu_X$.  If
$r(X)\ge1$ then, as $\nu_X\ge\nu$, this is at most
$\nu-\nu_X\le0$.  If $r(X)<1$ then $r\nu_X\ge r\nu$ and the
excess is at most $\nu(1-r)$.  The remaining statements follow
from Theorem~\ref{thm:excursion}.
\end{proof}

\begin{remark}\label{rem:ratiowin}
This also explains why the spike term yielded to
Proposition~\ref{prop:spike} while the intermediate range did
not.  A class of modulus comparable to $X$ appears in $O(1)$
windows, so the window drifts sum to $D(Y)$ and form a budget
that Markov's inequality can exploit.  A class of intermediate
modulus appears in \emph{every} larger window, so no such budget
exists and no Markov bound is available.  Numerically, among
nine dyadic windows up to $2\cdot10^6$, a modulus of size
$10^6$ occurs in two of them and a modulus of size $10^3$ in all
nine.

Two observations narrow what must be shown.  First, an
excursion requires $\nu>0$: the excess is $\nu-s(X)$ and
$s(X)\ge0$, so a system whose survivors vanish can never
produce one.  This removes every dying system, and in particular
disposes of the alarming ratio of
Remark~\ref{rem:intermediate}: the prime sieve has $\nu=0$, so
its $e^{\gamma}$ discrepancy threatens a bound, never the
conjecture.  (At finite $X$ the quantity $\nu_X$ is still
$e^{-\gamma}/\log X$ and far from its limit, which is easy to
confuse with $\nu$ itself.)

Second, the deviation is one-sided.  For each fixed $N$ the
survivors of the whole system lie inside those of the first $N$
classes, whose window density tends to $\nu_N$; hence
$\limsup_Xs(X)\le\nu_N$ for every $N$, and letting $N\to\infty$,
$\limsup_Xs(X)\le\nu$.  Since $\nu_X\to\nu>0$, it follows that
$\limsup_Xr(X)\le1$, which is what makes the signed and the
one-sided logarithmic means in Proposition~\ref{prop:ratio}
equivalent.  Combining,
the entire remaining content of the problem is that
\[
  r(X)\to1\ \text{ in logarithmic mean, for systems with }\nu>0 .
\]
Numerically, on a system with $\nu$ known exactly the ratio sits
at $1$ to three decimals across every window, the largest value
being $1.000136$; by contrast the \emph{sub}-systems of
Remark~\ref{rem:intermediate} reach $1.3297$, so the
one-sidedness is a property of full systems and not of
sieving in general.

The surviving hypothesis $\nu>0$ is itself structural.

\begin{proposition}[Structure of the open case]\label{prop:goodcyl}
Suppose $\nu>0$ and let $\varepsilon>0$.  Then there is a
modulus $M$ such that the residues $c$ whose cylinders are at
least $(1-\varepsilon)$-full of survivors carry all but
$\varepsilon$ of the survivor measure, and for each such $c$ the
induced system at $c$ contains \emph{no} pairwise-coprime
subfamily of divergent drift.
\end{proposition}

\begin{proof}
The survivor set $\Zhat\setminus\cU$ is closed of measure
$\nu>0$, so by the Lebesgue density theorem in $\Zhat$ almost
every one of its points has cylinders about it of survivor
proportion tending to $1$; choosing $M$ large enough along the
divisibility order gives the first assertion.  For the second,
fix such a $c$.  If the induced system at $c$ had a
pairwise-coprime subfamily of divergent drift, then by
Theorem~\ref{thm:D} its survivor measure would be $0$,
contradicting that the cylinder is $(1-\varepsilon)$-full.
\end{proof}

\begin{remark}
This says exactly why Theorem~\ref{thm:D} cannot reach the open
case: the hypothesis $\nu>0$ forbids, on the very cylinders
carrying the survivors, the coprime divergence that theorem
needs.  The separation is sharp in practice.  For the system with moduli $p^2$, where
$\nu=\prod_p(1-p^{-2})=6/\pi^2=0.6079\ldots$, the induced coprime
drift on every cylinder modulo $6$ is finite, being at most
$\sum_p6/p^2$ plus the contributions of $2$ and $3$; for the
system of all moduli, where $\nu=0$ whatever the residues
(Theorem~\ref{thm:D}, the primes being among the moduli), it
diverges on every cylinder.
\end{remark}

This suggests a question.  Since $\nu>0$
forbids a pairwise-coprime subfamily of divergent drift
(Theorem~\ref{thm:D}), the moduli must share primes heavily; and
if every modulus were divisible by one of finitely many primes,
a coprime subfamily could have at most that many members.  The two
simplest systems with $\nu>0$ and divergent drift -- the even
moduli at residue $0$ and the doubled construction -- are indeed
covered by the single prime $2$.

The naive reduction -- that divergent drift with $\nu>0$ should
force the moduli to be covered, up to a subfamily of convergent
drift, by finitely many primes, or the residues to be aligned --
fails: the many-protector system of the remark following
Theorem~\ref{thm:exhaust} has $\nu>0$, divergent drift, no finite
prime cover and misaligned residues, and its density exists by
cylinder exhaustion.  The question that would close the problem
is therefore the following.

\medskip\noindent\textbf{Question (refined).}
\emph{Can a system with divergent drift and $\nu>0$ pass the
checklist below in full -- that is, satisfy the hypotheses of
none of the theorems it names, nor reduce by finitely many
applications of Theorem~\ref{thm:exhaust} to one that does?}

\medskip\noindent
One kind of counterexample would be an \emph{infinite tower}: a
system each of whose reductions has divergent drift, positive
survivor measure and misaligned residues, never terminating in a
covered one; another would admit no proper reduction at all and
fail every other hypothesis on the list.  Each reduction passes to the moduli of the protecting
cylinders, which are divisors of the previous level's moduli and
must be sparse enough that their union leaves survivors; whether
infinitely many such levels can each carry divergent drift is
the question.  

Three constraints on any such tower are immediate.

\begin{proposition}[Tower constraints]\label{prop:tower}
In a tower of proper reductions -- each level coarsening
infinitely many classes -- the following hold.
\begin{enumerate}
\item[(i)] A class of level-zero modulus $n_0$ can be coarsened
at most $\log_2 n_0$ times, since each proper coarsening passes
to a proper divisor and at least halves the modulus; so among
the classes coarsened at level $K$, only those of level-zero
modulus $n_0\ge2^K$ occur.  (A class left uncoarsened -- a
singleton group on its own ball -- persists unchanged, so this
constrains the coarsened classes, not the whole system.)
\item[(ii)] No level contains a modulus equal to $1$, for that
cylinder is all of $\Zhat$ and its exhaustion forces
$\nu=0$.
\item[(iii)] No non-top level consists of prime moduli, since a
prime has no proper divisor exceeding $1$; the classes that
persist through $K$ levels have moduli with at least $K$ prime
factors counted with multiplicity.
\end{enumerate}
\end{proposition}

\begin{proof}
Each is a direct reading of the definition of a proper
reduction: the reduced modulus is a proper divisor of every
modulus in its group.
\end{proof}

So an infinite tower needs, at every level, infinitely many
classes still being coarsened, hence level-zero classes with
unboundedly long divisor chains, each level of
divergent drift and positive survivor measure, and no level
reducing to primes.  Attempting to build one exhibits a
dichotomy.  Take moduli $2^{a}q$ over sparse primes $q$ and
unbounded $a$, with residues aligned within each $q$ (the only
arrangement leaving survivors); level $K$ has moduli
$2^{a-K}q$, drift $\sim\sum1/q$ independent of $K$, and
$\nu=\prod(1-1/2q)$.  If $\sum1/q<\infty$ then $\nu>0$ but
every level has convergent drift, so the tower terminates in
Theorem~\ref{thm:conv}.  If $\sum1/q=\infty$ then every level
has divergent drift but $\nu=0$, and Corollary~\ref{cor:nuzero}
settles it.  Over twelve sparse primes, $\nu=0.9535$
with drift $0.093$ at every level; over four hundred dense
primes, drift $2.33$ with $\nu$ at $0.29$, and $\nu\to0$ as
primes are added, by Theorem~\ref{thm:D}.  The two conditions were not obtainable
together.

We offer the reason as a heuristic, not a proof: coarsening
divides out the structure the classes share, so the moduli at
the top of a tower tend toward pairwise coprimality, and there
Theorem~\ref{thm:D} forces exactly this dichotomy -- divergent
drift kills the survivors, while positive survivor measure
forces the drift to converge.  If that heuristic could be made
rigorous, towers would always be finite and, by the refined
Question, the problem would close.

The heuristic predicts that coarsening raises coprimality.  Under
the canonical coarsening $n\mapsto n/p(n)$, dividing out the
smallest prime, two measures over five levels were computed: the
fraction of modulus pairs with gcd $1$, and the drift share of
the largest pairwise-coprime subfamily.  For the even moduli the
pair fraction rose from $0$ to $0.63$; for all integers from
$0.61$ to $0.64$; for products $2^a3^bq$ from $0$ to $0.74$.  A
family built to resist -- the multiples of $6$, sharing two
primes -- stayed at $0$ for one step, since removing $2$ leaves
$3$ in common, then rose to $0.59$: a system sharing $k$ primes
needs $k$ steps, which the heuristic accommodates.  The
coprime-subfamily share rose in every family, from between
$0.12$ and $0.35$ to between $0.45$ and $0.75$.  This is
empirical support for the direction of the heuristic and
nothing more; a proof must control the rate, since a tower
could in principle coarsen forever without ever becoming
coprime enough.

The dichotomy just described is false at a single level: the
semiprimes have divergent reciprocal sum, no common factor, and
every pairwise coprime subfamily convergent, since such a
subfamily is a matching on the primes and the best matching
$6,35,143,\dots$ converges.  Whatever closes the problem must
come from coarsening, not from the shape of one level.  The
following system has divergent drift, misaligned residues and
no proper exhaustion step, and Theorem~\ref{thm:quasi} settles
it.

\begin{proposition}[Irreducible semiprimes]\label{prop:irred}
For each prime $p$ let $M_p$ consist of the primes above $p$
taken in increasing order until their reciprocal sum $S_p$
reaches $1/2$, and take the classes $(pq,a_{pq})$, $q\in M_p$,
with arbitrary residues.  Then:
\begin{enumerate}
\item[(i)] every cylinder $(d,c)$ with $d\ge2$ containing two
or more classes has $d=p$ prime, its induced moduli are distinct
primes, its induced drift is bounded, and its induced survivor
measure is positive; so no group of two or more classes dies
in the sense of Theorem~\ref{thm:exhaust}, and the system admits
no proper exhaustion step, only the trivial ones in which a
class is its own group;
\item[(ii)] the drift diverges;
\item[(iii)] $\nu=0$ for every choice of residues.
\end{enumerate}
\end{proposition}

\begin{proof}
(i) A cylinder whose modulus has two prime factors contains only
the class with that modulus.  For $d=p$ the classes are
$(pq,\cdot)$ with $q\in M_p$ and $(p'p,\cdot)$ with $p\in
M_{p'}$; the induced moduli $q$ and $p'$ are distinct primes,
so the induced survivor measure is $\prod(1-1/\ell)$ over them,
positive when $\sum1/\ell$ converges.  That sum is at most
$S_p+\sum_{p':\,p\in M_{p'}}1/p'$; the first term is below
$1/2+1/p$.  For the second, $p\in M_{p'}$ holds iff
$\sum_{p'<q<p}1/q<1/2$, a condition that only becomes easier as
$p'$ increases, so the admissible $p'$ are the primes in an
interval $[p'_0,p)$, and their reciprocal sum is
$1/p'_0+\sum_{p'_0<q<p}1/q<1/p'_0+1/2\le1$.  (ii) The drift is $\sum_pS_p/p\ge\frac12\sum_p1/p$.
(iii) The smallest prime of $pq$ is $p$.  Writing $M_{p,c}$ for
the $q\in M_p$ with $a_{pq}\equiv c\pmod p$ and $F_{p,c}$ for
the event that some $q\in M_{p,c}$ has $x\equiv a_{pq}\pmod q$,
\[
\mu(\Kh_p)=\sum_c\frac1p\Bigl(1-\prod_{q\in M_{p,c}}(1-1/q)\Bigr)
\ \ge\ \frac1p\sum_c\bigl(1-e^{-S_{p,c}}\bigr)
\ \ge\ \frac{1-e^{-1/2}}{p},
\]
the last step because $(1-e^{-a})+(1-e^{-b})\ge1-e^{-a-b}$.  So
$\sum_p\mu(\Kh_p)=\infty$.  For $p\ne p'$,
$\mu(\Kh_p\cap \Kh_{p'})=\sum_{c,c'}\mu(F_{p,c}\cap F_{p',c'})/pp'$.
Say $p<p'$, and condition on $x\equiv c\pmod p$ and
$x\equiv c'\pmod{p'}$, an event of measure $1/pp'$.  Given it,
$F_{p,c}\cap F_{p',c'}$ is contained in the union of three
events.  $E_0$: $p'\in M_{p,c}$ and $a_{pp'}\equiv c'\pmod{p'}$,
so that the witness $q=p'$ for the $p$-group is supplied by the
conditioning itself; this can happen for only one pair $(c,c')$,
namely $c\equiv a_{pp'}\pmod p$, $c'\equiv a_{pp'}\pmod{p'}$, and
contributes at most $1/pp'$ in all.  $E_1$: a common witness
$q\in M_{p,c}\cap M_{p',c'}$, $q\ne p'$, with
$a_{pq}\equiv a_{p'q}\pmod q$; given the conditioning the
coordinate at $q$ is still uniform, so $E_1$ has measure at most
$\sum1/q$ over such $q$, and each $q\in M_p\cap M_{p'}$ belongs
to exactly one pair $(c,c')$, so the $E_1$ terms total at most
$\min(S_p,S_{p'})/pp'$.  $E_2$: two distinct witnesses
$q\in M_{p,c}\setminus\{p'\}$ and $q''\in M_{p',c'}$ with
$q\ne q''$; both lie outside $\{p,p'\}$ (as $M_{p'}$ lies above
$p'$), so given the conditioning their coordinates are
independent and uniform, and the union bound over pairs gives
measure at most $S_{p,c}S_{p',c'}$, whose sum over $(c,c')$ is
$S_pS_{p'}$.  Altogether
$\mu(\Kh_p\cap \Kh_{p'})\le(S_pS_{p'}+\min(S_p,S_{p'})+1)/pp'
\le(7/4+o(1))/pp'$, so the ratio of Theorem~\ref{thm:quasi} is
at most $(7/4)/(1-e^{-1/2})^2+o(1)<11.4$.
\end{proof}

With $p\le150$ the system has $7{,}174$ classes on $582$ primes
up to $4{,}253$; computed exactly, the ratio is $1.889$ at
$N=149$ and falling from $2.159$, the least value of
$p\,\mu(\Kh_p)$ is $0.4815$ against the bound $0.3935$, the
largest induced drift on any cylinder is $0.886$, and the least
induced survivor measure $0.316$.  Restricting to four sparse
values of $p$ makes $\sum\mu(\Kh_p)=0.061$, the theorem falls
silent as it must, and the survivors have measure at least
$0.939$; a Monte Carlo draw of $3{,}000$ points of $\Zhat$
returned $0.941$.  So irreducibility does not yield a
counterexample: a top of a tower with $\sum_p\mu(\Kh_p)=\infty$
must carry smallest-prime kills whose second-moment ratio, along
the exhaustion by smallest prime, is unbounded.  Such correlation does \emph{not} force a
dying cylinder.

\begin{proposition}[Irreducible with survivors]\label{prop:irred2}
Let $p_1<p_2<\cdots$ be primes with $\sum_j1/p_j<1$, let
$S_j\to\infty$ with $\sum_jS_j/p_j=\infty$ -- for instance $p_j$
the $\lceil4j^{3/2}\rceil$th prime and $S_j=\sqrt j$: then
$\sum_j1/p_j<0.24$, using $p_k>k\log k$ beyond the first six
terms, while $p_k\sim k\log k$ gives $\sqrt j/p_j\sim1/(6j\log j)$,
whose sum diverges -- and let
group $j$ consist of the classes $(p_jq,a)$ with
$a\equiv1\pmod{p_j}$ and $a$ arbitrary modulo $q$, over primes
$q$ outside $\{p_i\}$ taken in order until $\sum1/q$ reaches
$S_j$.  Then the drift diverges, $\nu\ge1-\sum_j1/p_j>0$, no
cylinder holding two or more classes has induced survivor
measure $0$ (for the cylinder $\Zhat$ itself this is $\nu>0$),
so Theorem~\ref{thm:exhaust} admits no proper
reduction of the system; and
$\dlog(A)=\nu$ by Theorem~\ref{thm:confined}.
\end{proposition}

\begin{proof}
The drift is $\sum_jS_j/p_j$.  All kills lie in $\bigcup_jC_j$
with $C_j=\{x\equiv1\ (p_j)\}$, giving the bound on $\nu$.  A
cylinder of modulus at least $2$ holding two classes has prime
modulus: for $(p_j,1)$
the induced classes have the distinct prime moduli $q$ of group
$j$, with $\sum1/q<S_j+1$; for $(q,c)$ they have the distinct
prime moduli $p_j$ of the groups using $q$ with $a\equiv c\pmod
q$, with sum below $1$; in both cases the induced survivor
measure is $\prod(1-1/\ell)>0$.  Cylinders of composite modulus
hold at most one class.  Finally $G_1\cup\dots\cup G_J$ has
drift $\sum_{j\le J}S_j/p_j<\infty$, so Theorem~\ref{thm:conv}
covers it, and $\sum_j1/p_j<\infty$.
\end{proof}

With six groups the system has $6{,}500$ classes; every
cylinder holding two or more classes has prime modulus and
distinct prime induced moduli, the least induced survivor
measure is $0.0637$, the largest induced drift $2.449=\sqrt6$,
$\sum_{j\le6}1/p_j=0.1999$, and $6{,}490$ classes are
misaligned, with misaligned parts growing without bound.
Replacing the protectors by the first primes, so that
$\sum1/p_j$ diverges, sends the survivors down through
$0.487,0.432,0.391,0.359$ as $5,10,20,40$ groups are added, as
the theorem's failure predicts.

So the tower picture is incomplete in the other direction too:
a top that no exhaustion touches can be covered by confinement,
and the infinite protector tree that evades both -- deep
cylinders of measure bounded away from $0$, drift divergent in
total yet convergent along almost every branch -- is covered by
Theorem~\ref{thm:tree}.  What a counterexample must now evade is
the coprime-layered structure itself.  By
Proposition~\ref{prop:chain} that structure is a genuine
restriction: the nodes must share factors pairwise, so the
systems it reaches are those whose classes hang off a single
divisor chain of cylinders, one prime at a time.  Two natural
hopes are therefore closed off.  Not every system
admits a coprime-layered refinement -- the protector systems and
the semiprimes do not, though other theorems cover them -- and
relaxing coprimality of the layers loses the dying-branch lemma
outright, by Remark~\ref{rem:circular}.  The shape one might then fear -- moduli
spread across pairwise coprime cylinders with kills neither
summable nor correlated -- is ruled out by
Theorem~\ref{thm:coprot}, where positive survivor measure forces
the confinement.  Nested protectors are the tree theorem's;
pairwise coprime protectors are this one's.  Protectors that share factors
without nesting are Theorem~\ref{thm:overlap}'s, and there the
dichotomy is sharp: the protecting primes either have divergent
reciprocal sum, and then almost every integer is caught, or
convergent, and then the protectors are summable.  All three
regimes -- nested, coprime, overlapping -- are therefore closed,
each under the hypotheses of its theorem: finitely many children
and coprime layers for the trees, distinct attached primes for
the coprime protectors, every pair of the prime set with kills
bounded below for the overlapping ones.
What is not closed is the passage from these regimes to an
arbitrary system.  Each theorem assumes a given organisation of
the classes into groups with a stated relation between their
moduli, and no result here produces such an organisation from
nothing.  That is
too generous to the organisation theorems, whose hypotheses are
strong: Theorem~\ref{thm:coprot} wants its attached primes
distinct across all classes, so even Proposition~\ref{prop:irred2},
which reuses them, is covered by confinement rather than by it;
and the semiprimes are covered by Theorem~\ref{thm:quasi} without
any organisation at all.  The final form is therefore the
refined Question itself, read as a checklist.
A counterexample to the existence of the logarithmic density
would have to satisfy all of the following at once.
\begin{enumerate}
\item positive survivor measure (Corollary~\ref{cor:nuzero});
\item divergent drift (Theorem~\ref{thm:conv});
\item no pairwise coprime subfamily of divergent drift
(Theorem~\ref{thm:D});
\item smallest-prime kills whose measures are summable, or whose
second-moment ratio along the exhaustion by smallest prime is
unbounded (Theorem~\ref{thm:quasi});
\item for every modulus $M$, some cylinder modulo $M$ whose
induced system has no pairwise coprime subfamily of divergent
drift (Corollary~\ref{cor:dense});
\item no representation as finitely many aligned channels
(Theorems~\ref{thm:transl}, \ref{thm:chan});
\item unbounded misaligned parts, not confined to classes of
convergent drift (Theorem~\ref{thm:bounded},
Corollary~\ref{cor:mixed});
\item for some $\varepsilon_0>0$ and every modulus $M$, at least
a proportion $\varepsilon_0$ of the cylinders modulo $M$ neither
die nor have bounded misaligned parts (Theorems~\ref{thm:cyl},
\ref{thm:cylapp});
\item no grouping into dying cylinders whose reduced system has
density equal to its survivor measure
(Theorem~\ref{thm:exhaust}); and the same for its saturated
form, since by Theorem~\ref{thm:sat} a saturated form with
$\dlog(A_{\mathcal{P}})=\nu$ would force $\dlog(A)=\nu$;
\item no grouping into cylinders of summable measure with
covered finite unions (Theorem~\ref{thm:confined});
\item no coprime-layered structure (Theorem~\ref{thm:tree});
\item no pairwise coprime protectors with distinct attached
primes (Theorem~\ref{thm:coprot});
\item no protectors running over every pair of a set of primes,
carrying coprime layers of distinct attached primes with kills
bounded below (Theorem~\ref{thm:overlap}).
\end{enumerate}
Every system known to us fails this checklist at some line.
Whether it can be passed in full is the problem.

Theorem~\ref{thm:sat} gives a sufficient form without congruences:
a positive answer to the following question for every closed set
would settle the problem, though a negative one need not refute
it.
Let $\Ahat\subseteq\Zhat$ be closed, of positive Haar measure,
and let $\mathcal{H}$ be the cylinders meeting $\Ahat$ in a
Haar-null set that are maximal under inclusion, of divergent
total measure.  For an integer $x$, say $x$ is
\emph{caught} if $x\ge d$ for some $c+d\mathbb{Z}\in\mathcal{H}$
containing it.  Is the logarithmic density of the integers not
caught equal to $\mu(\Ahat)$?  Corollary~\ref{cor:anchor} says the
holes escaping any finite set of primes are either summable or
concentrated; the theorems above settle the cases in which the holes are
nested, pairwise coprime, or pairwise products, each under the
hypotheses of Theorems~\ref{thm:tree}, \ref{thm:coprot}
and~\ref{thm:overlap}.  What remains is a closed set whose holes
are concentrated in none of those ways, and whether such a set
exists is the form in which we leave the problem.

The role of the common factor can be isolated exactly.  Take the
odd moduli $3,5,7,\dots$ with residues in general position:
unbounded prime support, greatest common divisor $1$, drift
$\sim\tfrac12\log X$, and no alignment.  Its survivors fall, not quite monotonically, to
$0.0033$ at $X=4\cdot10^6$.  Now double every modulus and force
the residues even, changing nothing except that every kill is
confined to the even numbers: the survivors sit at $0.5016$, the
odd numbers being structurally unreachable.  The two systems
differ only by the protector, and only the protected one retains
survivors.  In its first form, before refinement, the question asked whether positive survivor
measure against divergent misaligned drift always requires such
a protector, hence a common factor; the many-protector system
answers that in the negative, and the refined Question asks
instead for a finite reduction.

We note that the block construction of Section~\ref{sec:routes} cannot be used
for this test: its drift grows like $s\log\log X$ and, summed over the eight
blocks $q\le23$, reaches only $0.05$, so neither the protected nor the
unprotected version has killed anything at reachable scales.  In the two simplest systems with $\nu>0$ and divergent drift,
the even moduli at residue $0$ and the doubled construction of
Section~\ref{sec:routes}, the survivors come from a cylinder that carries no
classes at all -- the odd numbers -- rather than from any
cancellation among the classes
acting there.  We know of no example of divergent drift in which a positive
survivor measure arises from genuine cancellation on a cylinder
whose induced drift diverges.  Whether such systems exist is itself
open.

What Proposition~\ref{prop:ratio} buys is not a proof but a
translation: the open statement is now that a Buchstab-type
ratio tends to $1$ on average, for arbitrary congruence systems
rather than for sieves by primes.  The classical theory
controls that ratio sharply in the prime case; whether it can be
controlled here we do not know.
\end{remark}

\begin{remark}
Theorem~\ref{thm:excursion} is an equivalence, not merely a
sufficient condition, and it relocates the difficulty: from controlling an infinite tail
to bounding the measure of the set of \emph{scales} at which
fresh classes over-kill.  By
Proposition~\ref{prop:driftinv} that control cannot come from
an averaged count of drift, so it must come from the arithmetic
of the residues.

Heuristically, to keep $w$ elevated on a set of
scales of positive logarithmic density one needs classes of
modulus comparable to $X$ at a positive proportion of scales,
hence essentially every integer as a modulus -- and such a
system kills everything, forcing $\nu=0$ and removing the
excursion.  This is the same self-defeating tension as in
Remark~\ref{rem:notvacuous} and in Section~\ref{sec:routes}, appearing here a
third time.  We have not turned it into a proof.
\end{remark}

\begin{remark}
Theorem~\ref{thm:reduction}, by contrast, is stated as a
sufficient condition only.  Necessity is not established here and does not follow formally: for fixed $k$ one
has $\lim_{X}\bigl(F(X)-F_k(X)\bigr)=\nu_k-\nu$, which tends to
$0$ with $k$ regardless, so all of the content lies in the
supremum, and a transient excursion could in principle exceed
the limit even when $\dlog(A)$ exists.
\end{remark}

\begin{remark}\label{rem:notvacuous}
It is worth asking whether the hypothesis can be met at all, and
the obvious attempt to defeat it fails for an instructive
reason.  One would try to manufacture transient excursions with a
block of moduli in $[N,2N]$.  Aligned at $0$ such a block is
useless: the density of integers with a divisor in $[N,2N]$
tends to $0$ as $N\to\infty$, by a theorem of
Erd\H{o}s~\cite{Er36}, as Besicovitch's construction requires.
With residues in general position the block's balls are nearly
independent, and it kills a profinite measure $d_N$ that stays
near $1-\prod_{N\le n\le2N}(1-1/n)\to\tfrac12$: measured on
$\Zhat$ with random residues, $0.516$, $0.513$ and $0.507$ at
$N=20,40,80$, against $0.426$, $0.413$ and $0.392$ for the same
moduli aligned at $0$.  So if such a block lies beyond the head
its fresh contribution is of order $d_N\nu_k$, giving
\[
  F(X)-F_k(X)\ \gtrsim\ d_N\nu_k
  \Bigl(1-\frac{\log N}{\log X}\Bigr),
\]
which at $X=N^2$ is about $d_N\nu_k/2$, apparently independent
of $k$.  But the strategy is self-defeating.  To keep this
bounded away from $0$ for every $k$ one needs infinitely many
blocks with $d_N$ bounded below; if the blocks use disjoint
sets of primes, so that their kill events are independent, $J$
of them leave survivor measure $\prod(1-d_N)\approx2^{-J}$, so
$\nu=0$ and the conclusion of the theorem holds trivially, while
blocks correlated by a common protector fall under
Theorem~\ref{thm:exhaust}.  What is left is a sequence of blocks
with $d_{N_j}$ bounded below, correlated neither by
independence nor by a common protector, and with $\nu>0$; we
have no argument excluding it, and it is a special case of the
refined Question.  So the natural counterexample strategy does
not work as stated, but neither is it refuted, and we do not
claim this settles necessity.
\end{remark}

\section{Four refuted routes, and a bound on a fifth}\label{sec:routes}

We record, with constructions, four natural approaches that do
not work, and a bound on a fifth.  These appear to us more useful
to a reader attacking the problem than the positive results.
The fifth is an attempt to \emph{refute} the conjecture rather
than prove it.

\subsection*{Divergence need not lie in a single channel}
Split the primes into two sets $Q$ and $R$, each of divergent
reciprocal sum, say by alternating; partition $R$ into
consecutive blocks $P_q$, one per $q\in Q$, each the shortest
block of harmonic weight $S_q:=\sum_{p\in P_q}1/p$ at least $s$,
so that $s\le S_q<s+1/\min P_q<s+1$, the cutoff large enough
that every block holds at least $q$ primes; take as moduli the products $qp$ for $p\in P_q$, list each
block in increasing order and give the class of modulus $qp_j$
the residue $j\bmod q$, so that the residue groups are the
arithmetic progressions of positions modulo $q$.  Classes from different
blocks are coprime, since the block primes lie in $R$ and the
indexing primes in $Q$, hence always compatible; classes within
a block have $\gcd=q$ exactly.  Hence a channel contains at most one residue
group per block, and the maximal channel decomposes as
$\sum_q(\text{heaviest group in block }q)$.  Since the weights
$1/p$ decrease along a block, the heaviest group is the one
containing the first prime, and for a decreasing sequence the
sum of every $q$th term is at most
$1/\min P_q+\tfrac1q\sum_{p\in P_q}1/p$; dividing by $q$ and
summing,
\[
  \text{maximal channel}\ \le\
  \sum_q\frac{1}{q\,\min P_q}\ +\ (s+1)\,P(2),
\]
where $P(2)=\sum_p p^{-2}$ is the prime zeta value.  Both terms
converge, the first because consecutive blocks of harmonic
weight at least $s$ force, by Mertens' theorem, the least prime
of the $k$th block to exceed $\exp(e^{ks})$ up to a constant,
so that $\sum_q1/(q\min P_q)$ is dominated by a doubly
exponentially decaying series.  The
total drift, meanwhile, is $\sum_q S_q/q\ge\sum_q s/q$, which
diverges by Mertens.  (The cruder bound $s\,P(2)\approx0.45s$ alone is not correct:
the block at $q=2$ contributes about $S_2/2\ge s/2$ on its own.)

\subsection*{Bounded channels need not force death}
Double every modulus in the previous construction and force even
residues.  Every kill set then lies in the even numbers, so every
odd number survives and the survivor density is at least
$1/2$, while the drift still diverges and the channels remain
bounded.

\subsection*{The $\varepsilon$-split fails}
The same doubled construction shows that no finite family of
channels captures all but $\varepsilon$ of a divergent drift.

\subsection*{Heilbronn--Rohrbach fails for shifted classes}
The inequality $\delta\ge\prod(1-1/n_i)$ of Heilbronn and
Rohrbach \cite{Heilbronn,Rohrbach}, generalised by Behrend
\cite{Behrend} and available for sets of multiples, is false for
shifted systems: for moduli $2,4,8$ with
residues $0,1,3$ the survivor density is $0.125$ against a bound
of $0.328125$.  Enumerating every system of three classes whose
moduli are a $3$-element subset of $\{2,\dots,12\}$ and whose
residues range independently over all residues modulo those
moduli gives
$\sum_{\{a<b<c\}\subset\{2,\dots,12\}}abc=53{,}130$ systems;
each is periodic modulo $\mathrm{lcm}(a,b,c)$, so its survivor
density is an exact rational, and $23{,}034$ of them --
$43.4\%$ -- fall strictly below the bound, the displayed system
being the extreme case.  The enumeration is specified here so
that it may be reproduced directly.

\subsection*{A bound for top-octave spikes}
Since $\overline{\delta}_{\log}(A)\le\nu$ is unconditional
(Theorem~\ref{thm:upper}), a counterexample would need the
killed set to exceed its own measure in upper \emph{logarithmic}
density, so $F$ must oscillate.  The natural attempt uses
classes of modulus comparable to $X$: a class with
$n\in(X/2,X]$ kills at most one element of that window, so $K$
of them can raise the window's killed density by $2K/X$ while
adding only $K/X$ to the drift.  Such spikes are real and large.
Taking the class $0\bmod n$ for every $n\in(Y/2,Y]$ and each
of $Y=2^{13}$, $2^{16}$, $2^{19}$, $2^{21}$ together, the killed
density of the dyadic window $(2^k,2^{k+1}]$ runs, for
$k=12,\dots,20$, through (the entries $1.000$ at $k=12,15,18,20$
being the four spiked windows)
\[
  1.000,\ 0.583,\ 0.527,\ 1.000,\ 0.692,\ 0.645,\ 1.000,
  \ 0.747,\ 1.000,
\]
an amplitude above $0.4$.  Yet $F$ does not oscillate at all:
evaluated at the tops $X=2^{k+1}$ of the same windows it rises
monotonically through
\[
0.077,\ 0.112,\ 0.140,\ 0.194,\ 0.222,\ 0.246,\ 0.286,\ 0.308,\ 0.341.
\]

The reason is structural.  A spike occupies one dyadic window,
of logarithmic length $\log 2$, so its excess above the
persistent level contributes $O(1/\log X)$ to the logarithmic
average and vanishes.  To move $F$ one would need the elevation
sustained over a constant fraction of $\log X$; but a class of
modulus $n$ kills density $1/n$ at \emph{every} scale above
$n$, so a sustained elevation is permanent rather than
transient, accumulates across scales, and drives the survivor
density to $0$ -- whereupon $\nu=0$ and the logarithmic density
exists trivially.  We state this last step as a heuristic, not a
theorem: no result here quantifies how successive excursions of
intermediate moduli interact.  What is proved is
Proposition~\ref{prop:spike}: isolated spikes cannot move the
logarithmic average when $h=0$.

\appendix
\section{Verification}\label{sec:verification}

The numerical assertions above were checked by computation,
deterministic except where a Monte Carlo value is quoted, in
which case the pseudo-random seed is fixed in the program;
several checks are accompanied by a control designed to fail if
the property being checked were absent.  The values on which the
text's explicit assertions rest are reproduced by a single
claims script; the Monte Carlo and window measurements quoted
as illustration are reproduced by the individual programs named
in the artifact, which is public at
\url{https://github.com/ibrahimmian36/Triarius}.  No theorem in this paper depends on these computations:
each is proved above, and the computations serve only to exhibit
the explicit examples and counterexamples, and in
Remark~\ref{rem:gap} to show why two plausible routes fail.  The programs, their recorded outputs and this paper form the
public, versioned artifact at that address.

\medskip\noindent\textbf{Provenance of the illustrative measurements.}
Each measurement quoted in a remark is produced by one named
program in the artifact's \texttt{probes} directory, run without
arguments: the many-protector and two-level-tower measurements
of Section~\ref{sec:reductions} by \path{verify_multiprotector.py} and
\path{verify_tower.py}; the window ratios and the
intermediate-range ratios by \path{verify_buchstab.py} and
\path{verify_overkill.py}; the induced-drift measurements of
the refinement remark by \path{verify_drift_invariance.py};
the coarsening statistics by \path{verify_coarsening.py}; the
semiprime, protector, tree and overlapping-protector instances
by \path{verify_quasi.py}, \path{verify_confined.py},
\path{verify_coprime_protectors.py}, \path{verify_tree.py}
and \path{verify_overlap.py}; the node assignments of the
chain constraint by \path{verify_organizable.py}; the
saturation reduction by \path{verify_saturation.py}; the
spike inequalities and the small-window excursions by
\path{verify_spike_bound.py} and \path{verify_excursion2.py}; the
protector test and the tower dichotomy of
Section~\ref{sec:reduction} by \path{verify_protector2.py} and
\path{verify_infinite_tower.py}; the window ratios and the
window occurrence counts of Remark~\ref{rem:ratiowin} by
\path{verify_narrow.py} and \path{verify_ratio.py}; the
Heilbronn--Rohrbach enumeration by \path{verify_paper.py}; the
spike display of Section~\ref{sec:routes} by the claims script;
and the catalogue audit of the eleven exhibited systems by
\path{verify_coverage.py}.  The
values that enter hypotheses or explicit assertions are
re-derived by \path{tools/verify_paper_claims.py} in one run,
whose output is archived with the artifact.  The disjointness bound
of Lemma~\ref{lem:A} was confirmed against an exhaustive
branch-and-bound maximum over all $104$ classes of modulus at most $14$, returning exactly $1.000000$.  The independence
identity behind Theorem~\ref{thm:D} was checked against a direct
sieve, and a non-coprime control deviates from it.  The
cylinder reduction of Theorem~\ref{thm:bounded} was verified
exactly on all non-empty cylinder-class pairs for a system with
misaligned parts $3,5,7$ and $M=105$.  The quantities in the proof of
Proposition~\ref{prop:irred} -- the reciprocal sums $S_p$, the
measures $\mu(\Kh_p)$ and $\mu(\Kh_p\cap \Kh_{p'})$ by
inclusion--exclusion over the independent prime coordinates,
and the induced drift and survivor measure of every cylinder --
were computed exactly for $p\le150$, with the sparse-$p$
restriction as the control on which the theorem's hypothesis
fails.  For Proposition~\ref{prop:irred2} the induced survivor
measure of every multi-class cylinder, the drift, and the
misalignment were computed exactly on the six-group truncation,
with the divergent-protector replacement as the control.  The
node assignments of Proposition~\ref{prop:chain} were searched
exhaustively for four truncated systems, confirming that a
depth-two slice of the tree and the multiples of $6$ are
coprime-layered and that the protector and semiprime systems
are not.  For
Theorem~\ref{thm:tree} an explicit tree of depth four --
children $k^2$ of $k^2+1$ residues at depth $k$, attached primes
from pools disjoint by depth and coprime to $2,5,17$, nodes of large drift in proportion $1/k^2$ -- has $617$ nodes and $23{,}361$
classes; its survivor measure $0.96432$, computed as the product
$\prod_q(1-k_q/q)$ from the proof of Lemma~\ref{lem:dying}, agrees with a Monte Carlo
draw of $20{,}000$ points ($0.9646$); at each of four levels
$N$ the first crossings of Lemma~\ref{lem:light} are pairwise
disjoint with total measure equal to $\mu\{g>N\}$ exactly, and the light drift is below $N$; and the window
survivor density agrees with the visible profinite measure to
three decimals at $X=10^4,10^5,10^6$ and to within $10^{-3}$ at
$4\times10^6$, which
illustrates the theorem and proves nothing.  Finally, as a safeguard against
misattribution, the catalogue of systems exhibited in this paper -- eleven of them, from the
aligned residues of Theorem~\ref{thm:transl} to the pairwise
products of Theorem~\ref{thm:overlap} -- was audited
mechanically: for each, the finitely checkable hypotheses of
the theorem the text cites for it are verified on a truncation,
the infinitary ones -- divergence, distinctness, escape -- being
established in the text where each system is introduced; the
finite checks pass.  The saturation reduction of
Theorem~\ref{thm:sat} was checked exactly on the finite systems
described there.

\medskip\noindent\textbf{Declaration of generative AI use.}
Large language models -- Anthropic's Claude (most recently Claude
Fable 5.1) and OpenAI's GPT-5.6 -- were used throughout this work:
to explore and draft arguments, to draft and revise the text, to
write the programs in the artifact, and, in separate sessions
without access to earlier reviews, to review the manuscript for
errors.  The authors directed the work, checked every proof, and
take full responsibility for the content.  Every number quoted in the paper is reproduced by a named
program of the artifact, whose recorded outputs it contains.

\begin{thebibliography}{9}

\bibitem{Behrend} F.~Behrend,
\emph{Generalization of an inequality of Heilbronn and
Rohrbach}, Bull. Amer. Math. Soc. \textbf{54} (1948), 681--684.

\bibitem{DE1936} H.~Davenport and P.~Erd\H{o}s,
\emph{On sequences of positive integers},
Acta Arith. \textbf{2} (1936), 147--151.

\bibitem{DE1951} H.~Davenport and P.~Erd\H{o}s,
\emph{On sequences of positive integers},
J. Indian Math. Soc. \textbf{15} (1951), 19--24.

\bibitem{Bloom} T.~F.~Bloom, \emph{Erd\H{o}s Problem \#25},
\url{https://www.erdosproblems.com/25}, accessed 4 September 2026.

\bibitem{Choj26} P.~Chojecki, \emph{Truncated congruence sieves
and Erd\H{o}s Problem 25}, note, 19 March 2026,
\url{https://www.ulam.ai/research/erdos25.pdf}.

\bibitem{HaRo66} H.~Halberstam and K.~F.~Roth, \emph{Sequences},
Vol.~I, Clarendon Press, Oxford, 1966.

\bibitem{TaoRogers} T.~Tao, \emph{Rogers' theorem on sieving},
blog post, 19 January 2026,
\url{https://terrytao.wordpress.com/2026/01/19/rogers-theorem-on-sieving/}.

\bibitem{TaoThread} T.~Tao, comment of 18 January 2026 quoting
G.~Tenenbaum with permission,
\url{https://www.erdosproblems.com/forum/thread/281#post-3370}.

\bibitem{Williams} D.~Williams,
\emph{Probability with Martingales},
Cambridge University Press, 1991.

\bibitem{KS} S.~Kochen and C.~Stone,
\emph{A note on the Borel--Cantelli lemma},
Illinois J. Math. \textbf{8} (1964), 248--251.

\bibitem{Er36} P.~Erd\H{o}s, \emph{A generalization of a theorem of
Besicovitch}, J. London Math. Soc. \textbf{11} (1936), 92--98.

\bibitem{Er48} P.~Erd\H{o}s, \emph{On the density of some
sequences of integers}, Bull. Amer. Math. Soc. \textbf{54} (1948),
685--692.

\bibitem{FFKPY} M.~Filaseta, K.~Ford, S.~Konyagin, C.~Pomerance
and G.~Yu, \emph{Sieving by large integers and covering systems of
congruences}, J. Amer. Math. Soc. \textbf{20} (2007), 495--517.

\bibitem{Hough} R.~Hough, \emph{Solution of the minimum modulus
problem for covering systems}, Ann. of Math. (2) \textbf{181}
(2015), 361--382.

\bibitem{HT92} R.~R.~Hall and G.~Tenenbaum, \emph{On Behrend
sequences}, Math. Proc. Cambridge Philos. Soc. \textbf{112}
(1992), 467--482.

\bibitem{Ten96} G.~Tenenbaum, \emph{On block Behrend sequences},
Math. Proc. Cambridge Philos. Soc. \textbf{120} (1996), 355--367.

\bibitem{Hall} R.~R.~Hall,
\emph{Sets of Multiples},
Cambridge Tracts in Mathematics \textbf{118},
Cambridge University Press, 1996.

\bibitem{Heilbronn} H.~Heilbronn,
\emph{On an inequality in the elementary theory of numbers},
Proc. Cambridge Philos. Soc. \textbf{33} (1937), 207--209.

\bibitem{Rohrbach} H.~Rohrbach,
\emph{Beweis einer zahlentheoretischen Ungleichung},
J. Reine Angew. Math. \textbf{177} (1937), 193--196.

\bibitem{Tenenbaum} G.~Tenenbaum,
\emph{Introduction to Analytic and Probabilistic Number
Theory},
Cambridge Studies in Advanced Mathematics \textbf{46},
Cambridge University Press, 1995.

\end{thebibliography}

\end{document}
